%
% $Header: /home/cvs/root/haskell-report/report/decls.verb,v 1.19 2003/01/13 13:08:55 simonpj Exp $
%
%**<title>The Haskell 98 Report: Declarations</title>
%*section 4
%**~header
\chapter{Declarations and Bindings}
\index{declaration}
\index{binding}
\label{declarations}

In this chapter, we describe the syntax and informal semantics of \Haskell{}
{\em declarations}.
% including their translations into
% the \Haskell{} kernel where appropriate.
% (see Chapter~\ref{formal-semantics} for a complete formal semantics).

\begin{flushleft}\it\begin{tabbing}
\hspace{0.6in}\=\hspace{3.1in}\=\kill
$\it module$\>\makebox[3.5em]{$\rightarrow$}$\it \makebox{\tt module}\ modid\ [exports]\ \makebox{\tt where}\ body$\\ 
$\it $\>\makebox[3.5em]{$|$}$\it body$\\ 
$\it body$\>\makebox[3.5em]{$\rightarrow$}$\it \makebox{\tt {\char'173}}\ impdecls\ \makebox{\tt ;}\ topdecls\ \makebox{\tt {\char'175}}$\\ 
$\it $\>\makebox[3.5em]{$|$}$\it \makebox{\tt {\char'173}}\ impdecls\ \makebox{\tt {\char'175}}$\\ 
$\it $\>\makebox[3.5em]{$|$}$\it \makebox{\tt {\char'173}}\ topdecls\ \makebox{\tt {\char'175}}$\\ 
$\it $\\ 
$\it topdecls$\>\makebox[3.5em]{$\rightarrow$}$\it topdecl_1\ \makebox{\tt ;}\ \ldots \ \makebox{\tt ;}\ topdecl_n$\>\makebox[3em]{}$\it (n\geq 1)$\\ 
$\it topdecl$\>\makebox[3.5em]{$\rightarrow$}$\it \makebox{\tt type}\ simpletype\ \makebox{\tt =}\ type$\\ 
$\it $\>\makebox[3.5em]{$|$}$\it \makebox{\tt data}\ [context\ \makebox{\tt =>}]\ simpletype\ \hprime{[}\makebox{\tt =}\ constrs\hprime{]}\ [deriving]$\\ 
$\it $\>\makebox[3.5em]{$|$}$\it \makebox{\tt newtype}\ [context\ \makebox{\tt =>}]\ simpletype\ \makebox{\tt =}\ newconstr\ [deriving]$\\ 
$\it $\>\makebox[3.5em]{$|$}$\it \makebox{\tt class}\ [scontext\ \makebox{\tt =>}]\ tycls\ tyvar\ [\makebox{\tt where}\ cdecls]$\\ 
$\it $\>\makebox[3.5em]{$|$}$\it \makebox{\tt instance}\ [scontext\ \makebox{\tt =>}]\ qtycls\ inst\ [\makebox{\tt where}\ idecls]$\\ 
$\it $\>\makebox[3.5em]{$|$}$\it \makebox{\tt default}\ \makebox{\tt (}type_1\ \makebox{\tt ,}\ \ldots \ \makebox{\tt ,}\ type_n\makebox{\tt )}$\>\makebox[3em]{}$\it \qquad\ (n\geq 0)$\\ 
$\it $\>\makebox[3.5em]{$|$}$\it \makebox{\tt foreign}\ fdecl$\\ 
$\it $\>\makebox[3.5em]{$|$}$\it decl$\\ 
$\it $\\ 
$\it decls$\>\makebox[3.5em]{$\rightarrow$}$\it \makebox{\tt {\char'173}}\ decl_1\ \makebox{\tt ;}\ \ldots \ \makebox{\tt ;}\ decl_n\ \makebox{\tt {\char'175}}$\>\makebox[3em]{}$\it (n\geq 0)$\\ 
$\it decl$\>\makebox[3.5em]{$\rightarrow$}$\it gendecl$\\ 
$\it $\>\makebox[3.5em]{$|$}$\it (funlhs\ |\ \hprime{pat})\ rhs$\\ 
$\it $\\ 
$\it cdecls$\>\makebox[3.5em]{$\rightarrow$}$\it \makebox{\tt {\char'173}}\ cdecl_1\ \makebox{\tt ;}\ \ldots \ \makebox{\tt ;}\ cdecl_n\ \makebox{\tt {\char'175}}$\>\makebox[3em]{}$\it (n\geq 0)$\\ 
$\it cdecl$\>\makebox[3.5em]{$\rightarrow$}$\it gendecl$\\ 
$\it $\>\makebox[3.5em]{$|$}$\it (funlhs\ |\ var)\ rhs$\\ 
$\it $\\ 
$\it idecls$\>\makebox[3.5em]{$\rightarrow$}$\it \makebox{\tt {\char'173}}\ idecl_1\ \makebox{\tt ;}\ \ldots \ \makebox{\tt ;}\ idecl_n\ \makebox{\tt {\char'175}}$\>\makebox[3em]{}$\it (n\geq 0)$\\ 
$\it idecl$\>\makebox[3.5em]{$\rightarrow$}$\it (funlhs\ |\ var)\ rhs$\\ 
$\it $\>\makebox[3.5em]{$|$}$\it $\>\makebox[3em]{}$\it (\tr{empty})$\\ 
$\it $\\ 
$\it gendecl$\>\makebox[3.5em]{$\rightarrow$}$\it vars\ \makebox{\tt ::}\ [context\ \makebox{\tt =>}]\ type$\>\makebox[3em]{}$\it (\tr{type\ signature})$\\ 
$\it $\>\makebox[3.5em]{$|$}$\it fixity\ [integer]\ ops$\>\makebox[3em]{}$\it (\tr{fixity\ declaration})$\\ 
$\it $\>\makebox[3.5em]{$|$}$\it $\>\makebox[3em]{}$\it (\tr{empty\ declaration})$\\ 
$\it $\\ 
$\it ops$\>\makebox[3.5em]{$\rightarrow$}$\it op_1\ \makebox{\tt ,}\ \ldots \ \makebox{\tt ,}\ op_n$\>\makebox[3em]{}$\it (n\geq 1)$\\ 
$\it vars$\>\makebox[3.5em]{$\rightarrow$}$\it var_1\ \makebox{\tt ,}\ \ldots \ \makebox{\tt ,}\ var_n$\>\makebox[3em]{}$\it (n\geq 1)$\\ 
$\it fixity$\>\makebox[3.5em]{$\rightarrow$}$\it \makebox{\tt infixl}\ |\ \makebox{\tt infixr}\ |\ \makebox{\tt infix}$
\end{tabbing}\end{flushleft}
\indexsyn{vars}%
\indexsyn{fixity}%
\indexsyn{ops}%
\indexsyn{topdecls}%
\indexsyn{topdecl}%
\indexsyn{gendecl}%
\indexsyn{decls}%
\indexsyn{decl}%
\indexsyn{cdecls}%
\indexsyn{cdecl}%
\indexsyn{idecls}%
\indexsyn{idecl}%
\indexsyn{fdecl}%

The declarations in the syntactic category \mbox{$\it topdecls$} are only allowed
at the top level of a \Haskell{} module (see
Chapter~\ref{modules}), whereas \mbox{$\it decls$} may be used either at the top level or
in nested scopes (i.e.~those within a \mbox{\tt let} or \mbox{\tt where} construct).

For exposition, we divide the declarations into
three groups: user-defined datatypes, consisting of \mbox{\tt type}, \mbox{\tt newtype},
and \mbox{\tt data} 
declarations (Section~\ref{user-defined-datatypes}); type classes and
overloading, consisting of \mbox{\tt class}, \mbox{\tt instance}, and \mbox{\tt default}
declarations (Section~\ref{overloading}); and nested declarations,
consisting of value bindings, type signatures, and fixity declarations
(Section~\ref{nested}).

%The \mbox{\tt module} declaration, along with \mbox{\tt import} and
%infix declarations, is described in Section~\ref{modules}.

\Haskell{} has several primitive datatypes that are ``hard-wired''
(such as integers and floating-point numbers), but most ``built-in''
datatypes are defined with normal \Haskell{} code, using normal \mbox{\tt type}
and \mbox{\tt data} declarations. % (see Section~\ref{user-defined-datatypes}).
These ``built-in'' datatypes are described in detail in
Section~\ref{basic-types}.

\section{Overview of Types and Classes}
\label{types-overview}

\Haskell{} uses a traditional
Hindley-Milner\index{Hindley-Milner type system}
polymorphic type system to provide a static type semantics
\cite{damas-milner82,hindley69}, but the type system has been extended with
{\em type classes} (or just {\em
classes}\index{class}) that provide 
a structured way to introduce {\em overloaded} functions.\index{type
class} \index{constructor class} \index{overloaded functions}

A \mbox{\tt class} declaration (Section~\ref{class-decls}) introduces a new
{\em type class} and the overloaded operations that must be
supported by any type that is an instance of that class.  An
\mbox{\tt instance} declaration (Section~\ref{instance-decls}) declares that a
type is an {\em instance} of a class and includes
the definitions of the overloaded operations---called {\em
class methods}---instantiated on the named type.
\index{class method}

For example, suppose we wish to overload the operations \mbox{\tt (+)} and
\mbox{\tt negate} on types \mbox{\tt Int} and \mbox{\tt Float}.  We introduce a new
type class called \mbox{\tt Num}:\nopagebreak[4]
\par
{\small
\bprog
{\small\begin{verbatim}
  class Num a  where          -- simplified class declaration for Num
    (+)    :: a -> a -> a     -- (Num is defined in the Prelude)
    negate :: a -> a
\end{verbatim}}
\eprog
}
This declaration may be read ``a type \mbox{\tt a} is an instance of the class
\mbox{\tt Num} if there are class methods \mbox{\tt (+)} and \mbox{\tt negate}, of the
given types, defined on it.''

We may then declare \mbox{\tt Int} and \mbox{\tt Float} to be instances of this class:
\bprog
{\small\begin{verbatim}
  instance Num Int  where     -- simplified instance of Num Int
    x + y       =  addInt x y
    negate x    =  negateInt x
  
  instance Num Float  where   -- simplified instance of Num Float
    x + y       =  addFloat x y
    negate x    =  negateFloat x
\end{verbatim}}
\eprog
where \mbox{\tt addInt}, \mbox{\tt negateInt}, \mbox{\tt addFloat}, and \mbox{\tt negateFloat} are assumed
in this case to be primitive functions, but in general could be any
user-defined function.  The first declaration above may be read
``\mbox{\tt Int} is an instance of the class \mbox{\tt Num} as witnessed by these
definitions (i.e.~class methods)\index{class method} for \mbox{\tt (+)} and \mbox{\tt negate}.''

More examples of type classes can be found in
the papers by Jones \cite{jones:cclasses} or Wadler and Blott
\cite{wadler:classes}. 
The term `type class' was used to describe the original \Haskell{} 1.0
type system; `constructor class' was used to describe an extension to
the original type classes.  There is no longer any reason to use two
different terms: in this report, `type class' includes both the
original \Haskell{} type classes and the constructor classes
introduced by Jones.

\subsection{Kinds}

To ensure that they are valid, type expressions are classified
into different {\em kinds}, \index{kind} which take one of two possible
forms:\nopagebreak[4]
\begin{itemize}

\item The symbol $\ast$ represents the kind of all nullary type
constructors.

\item If $\kappa_1$ and $\kappa_2$ are kinds, then $\kappa_1\rightarrow\kappa_2$
is the kind of types that take a type of kind $\kappa_1$ and return
a type of kind $\kappa_2$.
\end{itemize}
Kind inference checks the validity of type expressions 
in a similar way that type inference checks the validity of value expressions.  
However, unlike types, kinds are entirely
implicit and are not a visible part of the language.  Kind inference is discussed
in Section~\ref{kindinference}.

\subsection{Syntax of Types}
\label{type-syntax}
\index{type}
\label{types}

\begin{flushleft}\it\begin{tabbing}
\hspace{0.6in}\=\hspace{3.1in}\=\kill
$\it type$\>\makebox[3.5em]{$\rightarrow$}$\it btype\ [\makebox{\tt ->}\ type]$\>\makebox[3em]{}$\it (\tr{function\ type})$\\ 
$\it $\\ 
$\it btype$\>\makebox[3.5em]{$\rightarrow$}$\it [btype]\ atype$\>\makebox[3em]{}$\it (\tr{type\ application})$\\ 
$\it $\\ 
$\it atype$\>\makebox[3.5em]{$\rightarrow$}$\it gtycon$\\ 
$\it $\>\makebox[3.5em]{$|$}$\it tyvar$\\ 
$\it $\>\makebox[3.5em]{$|$}$\it \makebox{\tt (}\ type_1\ \makebox{\tt ,}\ \ldots \ \makebox{\tt ,}\ type_k\ \makebox{\tt )}$\>\makebox[3em]{}$\it (\tr{tuple\ type},\ k\geq 2)$\\ 
$\it $\>\makebox[3.5em]{$|$}$\it \makebox{\tt [}\ type\ \makebox{\tt ]}$\>\makebox[3em]{}$\it (\tr{list\ type})$\\ 
$\it $\>\makebox[3.5em]{$|$}$\it \makebox{\tt (}\ type\ \makebox{\tt )}$\>\makebox[3em]{}$\it (\tr{parenthesised\ constructor})$\\ 
$\it $\\ 
$\it gtycon$\>\makebox[3.5em]{$\rightarrow$}$\it qtycon$\\ 
$\it $\>\makebox[3.5em]{$|$}$\it \makebox{\tt ()}$\>\makebox[3em]{}$\it (\tr{unit\ type})$\\ 
$\it $\>\makebox[3.5em]{$|$}$\it \makebox{\tt []}$\>\makebox[3em]{}$\it (\tr{list\ constructor})$\\ 
$\it $\>\makebox[3.5em]{$|$}$\it \makebox{\tt (->)}$\>\makebox[3em]{}$\it (\tr{function\ constructor})$\\ 
$\it $\>\makebox[3.5em]{$|$}$\it \makebox{\tt (,}\{\makebox{\tt ,}\}\makebox{\tt )}$\>\makebox[3em]{}$\it (\tr{tupling\ constructors})$
\end{tabbing}\end{flushleft}
\indexsyn{type}%
\indexsyn{atype}%
\indexsyn{btype}%
\indexsyn{gtycon}%

\noindent
The syntax for \Haskell{} type expressions
\index{type expression}
\index{constructor expression}
is given above.  Just as data values are built using data
constructors, type values are built from \mbox{$\it type\ constructors$}.  As with
data constructors, the names of type constructors start with uppercase
letters.
Unlike data constructors, infix type constructors are not allowed (other than \mbox{\tt (->)}).

The main forms of type expression are as follows:\nopagebreak[4]
\begin{enumerate}

\item Type variables, written as identifiers beginning with
      a lowercase letter.  The kind of a variable is determined implicitly
      by the context in which it appears.

\item Type constructors.  Most type constructors are written as an identifier
      beginning with an uppercase letter.  For example:\nopagebreak[4]
      \begin{itemize}
      \item \mbox{\tt Char}, \mbox{\tt Int}, \mbox{\tt Integer}, \mbox{\tt Float}, \mbox{\tt Double} and \mbox{\tt Bool} are
            type constants with kind $\ast$.
      \item \mbox{\tt Maybe} and \mbox{\tt IO} are unary type
            constructors, and treated as types with
            kind $\ast\rightarrow\ast$.
      \item The declarations \mbox{\tt data\ T\ ...} or \mbox{\tt newtype\ T\ ...} add the type
            constructor \mbox{\tt T} to
            the type vocabulary.  The kind of \mbox{\tt T} is determined by
            kind inference.
      \end{itemize}
      Special syntax is provided for certain built-in type constructors:\nopagebreak[4]
      \begin{itemize}
      \item The {\em trivial type}\index{trivial type} is written as \mbox{\tt ()} and
            has kind $\ast$.
            It denotes the ``nullary tuple'' type, and has exactly one value,
            also written \mbox{\tt ()} (see Sections~\ref{unit-expression}
            and~\ref{basic-trivial}).
      \item The {\em function type} is written as \mbox{\tt (->)} and has\index{function type}
            kind $\ast\rightarrow\ast\rightarrow\ast$.
      \item The {\em list type}  is written as \mbox{\tt []} and has kind
            $\ast\rightarrow\ast$.\index{list}
      \item The {\em tuple types} are written as \mbox{\tt (,)},\index{tuple}
            \mbox{\tt (,,)}, and so on.  Their kinds are
            $\ast\rightarrow\ast\rightarrow\ast$,
            $\ast\rightarrow\ast\rightarrow\ast\rightarrow\ast$,  and
            so on.
      \end{itemize}
      Use of the \mbox{\tt (->)} and \mbox{\tt []} constants is described in more detail below.

\item Type application.  If $t_1$ is a type of kind
      $\kappa_1\rightarrow\kappa_2$ and $t_2$ is a type of kind $\kappa_1$,
      then $t_1~t_2$ is a type expression of kind $\kappa_2$.

\item A {\em parenthesized type}, having form \mbox{$\it \makebox{\tt (}t\makebox{\tt )}$}, is identical
      to the type \mbox{$\it t$}.

\end{enumerate}
For example, the type expression \mbox{\tt IO\ a} can be understood as the application
of a constant, \mbox{\tt IO}, to the variable \mbox{\tt a}.  Since the \mbox{\tt IO} type
constructor has kind 
$\ast\rightarrow\ast$, it follows that both the variable \mbox{\tt a} and the whole
expression, \mbox{\tt IO\ a}, must have kind $\ast$.
In general, a process of {\em kind inference}\index{kind}\index{kind inference}
(see Section~\ref{kindinference})
is needed to determine appropriate kinds for user-defined datatypes, type
synonyms, and classes.

Special syntax is provided to allow certain type expressions to be written
in a more traditional style:\nopagebreak[4]
\begin{enumerate}

\item A {\em function type}\index{function type} has the form
\mbox{$\it t_1\ \makebox{\tt ->}\ t_2$}, which is equivalent to the type
\mbox{$\it \makebox{\tt (->)}\ t_1\ t_2$}.  Function arrows associate to the right.
For example, \mbox{\tt Int\ ->\ Int\ ->\ Float} means \mbox{\tt Int\ ->\ (Int\ ->\ Float)}.

\item A {\em tuple type}\index{tuple type} has the form 
\mbox{$\it \makebox{\tt (}t_1\makebox{\tt ,}\ \ldots \ \makebox{\tt ,}\ t_k\makebox{\tt )}$} where \mbox{$\it k\geq 2$}, which is equivalent to
the type \mbox{$\it \makebox{\tt (,}$\ldots$\makebox{\tt ,)}\ t_1\ \ldots \ t_k$} where there are
$k-1$ commas between the parenthesis.  It denotes the
type of \mbox{$\it k$}-tuples with the first component of type \mbox{$\it t_1$}, the second
component of type \mbox{$\it t_2$}, and so on (see Sections~\ref{tuples}
and \ref{basic-tuples}).

\item A {\em list type}\index{list type} has the form \mbox{$\it \makebox{\tt [}t\makebox{\tt ]}$},
which is equivalent to the type \mbox{$\it \makebox{\tt []}\ t$}.
It denotes the type of lists with elements of type \mbox{$\it t$} (see
Sections~\ref{lists} and \ref{basic-lists}).

\end{enumerate}
These special syntactic forms always denote the built-in type constructors
for functions, tuples, and lists, regardless of what is in scope.
In a similar way, the prefix type constructors \mbox{\tt (->)}, \mbox{\tt []}, \mbox{\tt ()}, \mbox{\tt (,)}, 
and so on, always denote the built-in type constructors; they 
cannot be qualified, nor mentioned in import or export lists (Chapter~\ref{modules}).
(Hence the special production, ``gtycon'', above.)

Although the list and tuple types have special syntax, their semantics 
is the same as the equivalent user-defined algebraic data types.

Notice that expressions and types have a consistent syntax.
If \mbox{$\it t_i$} is the type of
expression or pattern \mbox{$\it e_i$}, then the expressions \mbox{$\it \makebox{\tt ({\char'134}}\ e_1\ \makebox{\tt ->}\ e_2\makebox{\tt )}$},
\mbox{$\it \makebox{\tt [}e_1\makebox{\tt ]}$}, and \mbox{$\it \makebox{\tt (}e_1,e_2\makebox{\tt )}$} have the types \mbox{$\it \makebox{\tt (}t_1\ \makebox{\tt ->}\ t_2\makebox{\tt )}$},
\mbox{$\it \makebox{\tt [}t_1\makebox{\tt ]}$}, and \mbox{$\it \makebox{\tt (}t_1,t_2\makebox{\tt )}$}, respectively.

\index{quantification}
With one exception (that of the distinguished type variable
in a class declaration (Section~\ref{class-decls})), the
type variables in a \Haskell{} type expression
are all assumed to be universally quantified; there is no explicit
syntax for universal quantification~\cite{damas-milner82}.
% \cite{damas-milner82,reynolds90}.
For example, the type expression
\mbox{\tt a\ ->\ a} denotes the type \mbox{$\it \forall\ a.~a\ \rightarrow\ a$}.
For clarity, however, we often write quantification explicitly
when discussing the types of \Haskell{} programs.  When we write an
explicitly quantified type, the scope of the \mbox{$\it \forall$} extends as far
to the right as possible; for example, \mbox{$\it \forall\ a.~a\ \rightarrow\ a$} means
\mbox{$\it \forall\ a.~(a\ \rightarrow\ a)$}.

%Every type variable appearing in a signature
%is universally quantified over that signature.  This last
%constraint implies that signatures such as:
%\bprog
%@
%       \ x -> ([x] :: [a])
%@
%\eprog
%are not valid, because this declares \mbox{\tt [x]} to be of type 
%\mbox{$\it \forall\ a.\makebox{\tt [}a\makebox{\tt ]}$}.  In contrast, this {\em is} a valid signature:
%\mbox{\tt ({\char'134}\ x\ ->\ [x])\ ::\ a\ ->\ [a]}; it declares that \mbox{\tt ({\char'134}\ x\ ->\ [x])} has type
%\mbox{$\it \forall\ a.a\ \makebox{\tt ->}\ \makebox{\tt [}a\makebox{\tt ]}$}.

\subsection{Syntax of Class Assertions and Contexts}
\index{class assertion}
\index{context}
\label{classes&contexts}

\begin{flushleft}\it\begin{tabbing}
\hspace{0.6in}\=\hspace{3.1in}\=\kill
$\it context$\>\makebox[3.5em]{$\rightarrow$}$\it class$\\ 
$\it $\>\makebox[3.5em]{$|$}$\it \makebox{\tt (}\ class_1\ \makebox{\tt ,}\ \ldots \ \makebox{\tt ,}\ class_n\ \makebox{\tt )}$\>\makebox[3em]{}$\it (n\geq 0)$\\ 
$\it class$\>\makebox[3.5em]{$\rightarrow$}$\it qtycls\ tyvar$\\ 
$\it $\>\makebox[3.5em]{$|$}$\it qtycls\ \makebox{\tt (}\ tyvar\ atype_1\ \ldots \ atype_n\ \makebox{\tt )}$\>\makebox[3em]{}$\it (n\geq 1)$\\ 
$\it qtycls$\>\makebox[3.5em]{$\rightarrow$}$\it [\ modid\ \makebox{\tt .}\ ]\ tycls$\\ 
$\it tycls$\>\makebox[3.5em]{$\rightarrow$}$\it conid$\\ 
$\it tyvar$\>\makebox[3.5em]{$\rightarrow$}$\it varid$
\end{tabbing}\end{flushleft}
\indexsyn{context}%
\indexsyn{class}%
\indexsyn{simpleclass}%
\indexsyn{tycls}%
\indexsyn{tyvar}%

A {\em class assertion} has form \mbox{$\it qtycls\ tyvar$}, and
indicates the membership of the type \mbox{$\it tyvar$} in the class
\mbox{$\it qtycls$}.  A class identifier begins with an uppercase
letter.
A {\em context} consists of zero or more class assertions,
and has the general form
\[
\mbox{$\it \makebox{\tt (}\ C_1\ u_1,\ \ldots ,\ C_n\ u_n\ \makebox{\tt )}$}
\]
where \mbox{$\it C_1,\ \ldots ,\ C_n$} are class identifiers, and each of the \mbox{$\it u_1,\ \ldots ,\ u_n$} is
either a type variable, or the application of type variable to one or more types.
The outer parentheses may be omitted when \mbox{$\it n=1$}.  In
general, we use \mbox{$\it cx$} to denote a context and we write \mbox{$\it cx\ \makebox{\tt =>}\ t$} to
indicate the type \mbox{$\it t$} restricted by the context \mbox{$\it cx$}.
The context \mbox{$\it cx$} must only contain type variables referenced in \mbox{$\it t$}.
For convenience,
we write \mbox{$\it cx\ \makebox{\tt =>}\ t$} even if the context \mbox{$\it cx$} is empty, although in this
case the concrete syntax contains no \mbox{\tt =>}.

\subsection{Semantics of Types and Classes}
\label{type-semantics}

In this section, we provide informal details of the type system.
% the formal semantics is described in Chapter~\ref{static-semantics}
(Wadler and Blott \cite{wadler:classes} and Jones
\cite{jones:cclasses} discuss type
and constructor classes, respectively, in more detail.)

%A type is a {\em monotype\/}\index{monotype} if it contains no type
%variables, and is {\em monomorphic\/}
%\index{monomorphic type}
%if it contains type variables
%but is not polymorphic (in Milner's original terminology,
%it is monomorphic if it contains no generic type variables).

The \Haskell{} type system attributes a {\em type} to each
\index{type}
expression in the program.  In general, a type is of the form
\mbox{$\it \forall\ \overline{u}.~cx\ \Rightarrow\ t$},
where \mbox{$\it \overline{u}$} is a set of type variables \mbox{$\it u_1,\ \ldots ,\ u_n$}.
In any such type, any of the universally-quantified type variables \mbox{$\it u_i$}
that are free in \mbox{$\it cx$} must also be free in \mbox{$\it t$}.
Furthermore, the context \mbox{$\it cx$} must be of the form given above in
Section~\ref{classes&contexts}.  For example, here are some
valid types:
\bprog
{\small\begin{verbatim}
  Eq a => a -> a
  (Eq a, Show a, Eq b) => [a] -> [b] -> String
  (Eq (f a), Functor f) => (a -> b) -> f a -> f b -> Bool
\end{verbatim}}
\eprog
In the third type, the constraint \mbox{\tt Eq\ (f\ a)} cannot be made
simpler because \mbox{\tt f} is universally quantified.

The type of an expression \mbox{$\it e$} depends 
on a {\em type environment}\index{type environment} that gives types 
for the free variables in \mbox{$\it e$}, and a
{\em class environment}\index{class environment} that 
declares which types are instances of which classes (a type becomes
an instance of a class only via the presence of an
\mbox{\tt instance} declaration or a \mbox{\tt deriving} clause).

Types are related by a generalization preorder
\index{generalization preorder}
(specified below);
the most general type, up to the equivalence induced by the generalization preorder,
that can be assigned to a particular
expression (in a given environment) is called its {\em
principal type}.
\index{principal type}
\Haskell{}'s extended Hindley-Milner type system can infer the principal
type of all expressions, including the proper use of overloaded
class methods (although certain ambiguous overloadings could arise, as
described in Section~\ref{default-decls}).  Therefore, explicit typings (called
{\em type signatures})
\index{type signature}
are usually optional (see
Sections~\ref{expression-type-sigs} and~\ref{type-signatures}).

The type \mbox{$\it \forall\ \overline{u}.~cx_1\ \Rightarrow\ t_1$} is
{\em more general than} the 
type \mbox{$\it \forall\ \overline{w}.~cx_2\ \Rightarrow\ t_2$} if and only if there is 
a substitution \mbox{$\it S$} whose domain is \mbox{$\it \overline{u}$} such that:
\begin{itemize}
\item \mbox{$\it t_2$} is identical to \mbox{$\it S(t_1)$}.
\item Whenever \mbox{$\it cx_2$} holds in the class environment, \mbox{$\it S(cx_1)$} also holds.
\end{itemize}

A value of type
\mbox{$\it \forall\ \overline{u}.~cx\ \Rightarrow\ t$},
may be instantiated at types \mbox{$\it \overline{s}$} if and only if
the context \mbox{$\it cx[\overline{s}/\overline{u}]$} holds.
For example, consider the function \mbox{\tt double}:
\bprog
{\small\begin{verbatim}
  double x = x + x
\end{verbatim}}
\eprog
The most general type of \mbox{\tt double} is
\mbox{$\it \forall\ a.~\makebox{\tt Num}~a\ \Rightarrow\ a\ \rightarrow\ a$}.
\mbox{\tt double} may be applied to values of type \mbox{\tt Int} (instantiating \mbox{$\it a$} to
\mbox{\tt Int}), since \mbox{\tt Num\ Int} holds, because \mbox{\tt Int} is an instance of the class \mbox{\tt Num}.
However, \mbox{\tt double} may not normally be applied to values
of type \mbox{\tt Char}, because \mbox{\tt Char} is not normally an instance of class \mbox{\tt Num}.  The
user may choose to declare such an instance, in which case \mbox{\tt double} may
indeed be applied to a \mbox{\tt Char}.


\section{User-Defined Datatypes}
\index{datatype}
\label{user-defined-datatypes}

In this section, we describe algebraic datatypes (\mbox{\tt data}
declarations), renamed datatypes (\mbox{\tt newtype} declarations), and type
synonyms (\mbox{\tt type} declarations).  These declarations may only appear at
the top level of a module.

\subsection{Algebraic Datatype Declarations}
\index{algebraic datatype}
\label{datatype-decls}

\begin{flushleft}\it\begin{tabbing}
\hspace{0.6in}\=\hspace{3.1in}\=\kill
$\it topdecl$\>\makebox[3.5em]{$\rightarrow$}$\it \makebox{\tt data}\ [context\ \makebox{\tt =>}]\ simpletype\ \hprime{[}\makebox{\tt =}\ constrs\hprime{]}\ [deriving]$\\ 
$\it $\\ 
$\it simpletype$\>\makebox[3.5em]{$\rightarrow$}$\it tycon\ tyvar_1\ \ldots \ tyvar_k$\>\makebox[3em]{}$\it (k\geq 0)$\\ 
$\it $\\ 
$\it constrs$\>\makebox[3.5em]{$\rightarrow$}$\it constr_1\ \makebox{\tt |}\ \ldots \ \makebox{\tt |}\ constr_n$\>\makebox[3em]{}$\it (n\geq 1)$\\ 
$\it constr$\>\makebox[3.5em]{$\rightarrow$}$\it con\ [\makebox{\tt !}]\ atype_1\ \ldots \ [\makebox{\tt !}]\ atype_k$\>\makebox[3em]{}$\it (\tr{arity}\ con\ =\ k,\ k\geq 0)$\\ 
$\it $\>\makebox[3.5em]{$|$}$\it (btype\ |\ \makebox{\tt !}\ atype)\ conop\ (btype\ |\ \makebox{\tt !}\ atype)$\>\makebox[3em]{}$\it (\tr{infix}\ conop)$\\ 
$\it $\>\makebox[3.5em]{$|$}$\it con\ \makebox{\tt {\char'173}}\ fielddecl_1\ \makebox{\tt ,}\ \ldots \ \makebox{\tt ,}\ fielddecl_n\ \makebox{\tt {\char'175}}$\>\makebox[3em]{}$\it (n\geq 0)$\\ 
$\it fielddecl$\>\makebox[3.5em]{$\rightarrow$}$\it vars\ \makebox{\tt ::}\ (type\ |\ \makebox{\tt !}\ atype)$\\ 
$\it $\\ 
$\it deriving$\>\makebox[3.5em]{$\rightarrow$}$\it \makebox{\tt deriving}\ (dclass\ |\ \makebox{\tt (}dclass_1\makebox{\tt ,}\ \ldots \ \makebox{\tt ,}\ dclass_n\makebox{\tt )})$\>\makebox[3em]{}$\it (n\geq 0)$\\ 
$\it dclass$\>\makebox[3.5em]{$\rightarrow$}$\it qtycls$
\end{tabbing}\end{flushleft}
\index{topdecl@{\em topdecl} (\mbox{\tt data})}%
\indexsyn{simpletype}%
\indexsyn{constrs}%
\indexsyn{constr}%
\indexsyn{fielddecl}%
\indexsyn{deriving}%
\indexsyn{dclass}%
\index{data declaration@{\tt data} declaration}
The precedence\index{precedence} for \mbox{$\it constr$} is the same as that for
expressions---normal constructor application has higher precedence
than infix constructor application (thus \mbox{\tt a\ :\ Foo\ a} parses as 
\mbox{\tt a\ :\ (Foo\ a)}).

An algebraic datatype declaration has the form:
\[
\mbox{$\it \makebox{\tt data}\ cx\ \makebox{\tt =>}\ T\ u_1\ \ldots \ u_k\ \makebox{\tt =}\ K_1\ t_{11}\ \ldots \ t_{1k_1}\ \makebox{\tt |}\ \cdots\ \makebox{\tt |}\ K_n\ t_{n1}\ \ldots \ t_{nk_n}$}
\]
where \mbox{$\it cx$} is a context.
%\index{context!in data declaration@in {\tt data} declaration}
This declaration
introduces a new {\em type constructor} \mbox{$\it T$} with \hprime{zero} or more constituent {\em data
constructors} \mbox{$\it K_1,\ \ldots ,\ K_n$}.
\index{data constructor}\index{type constructor}
In this Report, the unqualified term ``constructor'' always means ``data constructor''.

The types of the data constructors are given by:
\[
\mbox{$\it K_i\ ::\ \forall\ u_1\ \ldots \ u_k.~\ cx_i\ \Rightarrow\ t_{i1}\ \rightarrow\ \cdots\ \rightarrow\ t_{ik_i}\ \rightarrow\ (T\ u_1\ \ldots \ u_k)$}
\]
where \mbox{$\it cx_i$} is the largest subset of \mbox{$\it cx$} that constrains only those type
variables free in the types \mbox{$\it t_{i1},\ \ldots ,\ t_{ik_i}$}.
The type variables \mbox{$\it u_1$} through \mbox{$\it u_k$} must be distinct and may appear
in \mbox{$\it cx$} and the \mbox{$\it t_{ij}$}; it is a static error
for any other type variable to appear in \mbox{$\it cx$} or on the right-hand-side.
The new type constant \mbox{$\it T$} has a kind of the form
$\kappa_1\rightarrow\ldots\rightarrow\kappa_k\rightarrow\ast$
where the kinds \mbox{$\it \kappa_i$} of the argument variables \mbox{$\it u_i$} are
determined by kind inference
as described in Section~\ref{kindinference}\index{kind}\index{kind inference}.
This means that \mbox{$\it T$} may be used in type expressions with anywhere
between \mbox{$\it 0$} and \mbox{$\it k$} arguments.

For example, the declaration
\bprog
{\small\begin{verbatim}
  data Eq a => Set a = NilSet | ConsSet a (Set a)
\end{verbatim}}
\eprog
introduces a type constructor \mbox{\tt Set} of kind $\ast\rightarrow\ast$, and constructors \mbox{\tt NilSet} and
\mbox{\tt ConsSet} with types
\[\ba{ll}
\mbox{\tt NilSet}  & \mbox{$\it ::\ \forall\ a.~\ \makebox{\tt Set}~\ a$} \\
\mbox{\tt ConsSet} & \mbox{$\it ::\ \forall\ a.~\ \makebox{\tt Eq}~\ a\ \Rightarrow\ a\ \rightarrow\ \makebox{\tt Set}~\ a\ \rightarrow\ \makebox{\tt Set}~\ a$}
\ea\]
In the example given, the overloaded
type for \mbox{\tt ConsSet} ensures that \mbox{\tt ConsSet} can only be applied to values whose
type is an instance of the class \mbox{\tt Eq}.  
Pattern matching against \mbox{\tt ConsSet} also gives rise to an \mbox{\tt Eq\ a} constraint. 
For example: 
\bprog
{\small\begin{verbatim}
  f (ConsSet a s) = a
\end{verbatim}}
\eprog
the function \mbox{\tt f} has inferred type \mbox{\tt Eq\ a\ =>\ Set\ a\ ->\ a}.
The context in the \mbox{\tt data} declaration has no other effect whatsoever.

The visibility of a datatype's constructors (i.e.~the ``abstractness''
\index{abstract datatype}
of the datatype) outside of the module in which the datatype is
defined is controlled by the form of the datatype's name in the export
list as described in Section~\ref{abstract-types}.

The optional \mbox{$\it \makebox{\tt deriving}$} part of a \mbox{\tt data} declaration has to do
with {\em derived instances}, and is described in
Section~\ref{derived-decls}.

\paragraph{Labelled Fields}
\label{field-labels}
\index{field label}
A data constructor of arity \mbox{$\it k$} creates an object with \mbox{$\it k$} components.
These components are normally accessed positionally as arguments to the
constructor in expressions or patterns.  For large datatypes it is
useful to assign \mbox{$\it field\ labels$} to the components of a data object.
This allows a specific field to be referenced independently of its
location within the constructor.

A constructor definition in a \mbox{\tt data} declaration may assign labels to the
fields of the constructor, using the record syntax (\mbox{\tt C\ {\char'173}\ ...\ {\char'175}}).
Constructors using field labels may be freely mixed with constructors
without them. 
A constructor with associated field labels may still be used as an
ordinary constructor; features using labels are
simply a shorthand for operations using an underlying positional
constructor.  The arguments to the positional constructor occur in the
same order as the labeled fields.  For example, the declaration
\bprog
{\small\begin{verbatim}
  data C = F { f1,f2 :: Int, f3 :: Bool }
\end{verbatim}}
\eprog
defines a type and constructor identical to the one produced by
\bprog
{\small\begin{verbatim}
  data C = F Int Int Bool
\end{verbatim}}
\eprog
Operations using field labels are described in Section~\ref{field-ops}.
A \mbox{\tt data} declaration may use the same field label in multiple
constructors as long as the typing of the field is the same in all
cases after type synonym expansion.  A label cannot be shared by
more than one type in scope.  Field names share the top level namespace
with ordinary variables and class methods and must not conflict with
other top level names in scope.

The pattern \mbox{\tt F\ {\char'173}{\char'175}} matches any value built with constructor \mbox{\tt F}, 
{\em whether or not \mbox{\tt F} was declared with record syntax}. 

\paragraph{Strictness Flags}
\label{strictness-flags}
\index{strictness flag}
Whenever a data constructor is applied, each argument to the
constructor is evaluated if and only if the corresponding type in the
algebraic datatype declaration has a strictness flag, denoted by
an exclamation point, ``\mbox{\tt !}''.
\index{"!@{\tt  {\char'041}}}
Lexically, ``\mbox{\tt !}'' is an ordinary varsym not a \mbox{$\it reservedop$}; 
it has special significance only in the context of the argument types of 
a data declaration.


\outline{
\paragraph*{Translation:}
A declaration of the form
\[
\mbox{$\it \makebox{\tt data}\ cx\ \makebox{\tt =>}\ T\ u_1\ \ldots \ u_k\ \makebox{\tt =}\ \ldots \ \makebox{\tt |}\ K\ s_{1}\ \ldots \ s_{n}\ \makebox{\tt |}\ \ldots $}
\]
where each \mbox{$\it s_i$} is either of the form \mbox{\tt !}\mbox{$\it t_i$} or \mbox{$\it t_i$}, replaces
every occurrence of \mbox{$\it K$} in an expression by 
\[
\mbox{$\it \makebox{\tt ({\char'134}\ }x_1\ \ldots \ x_n\ \makebox{\tt ->}\ (\ ((K\ op_1\ x_1)\ op_2\ x_2)\ \ldots \ )\ op_n\ x_n)$}
\]
where \mbox{$\it op_i$} is the non-strict apply function \mbox{\tt {\char'44}} if \mbox{$\it s_i$} is of the form \mbox{$\it t_i$},
and \mbox{$\it op_i$} is the strict apply function \mbox{\tt {\char'44}!} (see
Section~\ref{strict-eval}) if \mbox{$\it s_i$} is of the form \mbox{$\it \makebox{\tt !}\ t_i$}.  Pattern
matching on \mbox{$\it K$} is not affected by strictness flags.
}


\subsection{Type Synonym Declarations}
\index{type synonym}
\label{type-synonym-decls}

\begin{flushleft}\it\begin{tabbing}
\hspace{0.6in}\=\hspace{3.1in}\=\kill
$\it topdecl$\>\makebox[3.5em]{$\rightarrow$}$\it \makebox{\tt type}\ simpletype\ \makebox{\tt =}\ type$\\ 
$\it simpletype$\>\makebox[3.5em]{$\rightarrow$}$\it tycon\ tyvar_1\ \ldots \ tyvar_k$\>\makebox[3em]{}$\it (k\geq 0)$
\end{tabbing}\end{flushleft}
\index{topdecl@{\em topdecl} (\mbox{\tt type})}%
\indexsyn{simpletype}%
A type synonym declaration introduces a new type that
is equivalent to an old type.  It has the form
\[
\mbox{$\it \makebox{\tt type}\ T\ u_1\ \ldots \ u_k\ \makebox{\tt =}\ t$}
\]
which introduces a new type constructor, \mbox{$\it T$}.  The type \mbox{$\it (T\ t_1\ \ldots \\
\it t_k)$} is equivalent to the type \mbox{$\it t[t_1/u_1,\ \ldots ,\ t_k/u_k]$}.  The type
variables \mbox{$\it u_1$} through \mbox{$\it u_k$} must be distinct and are scoped only
over \mbox{$\it t$}; it is a static error for any other type variable to appear
in \mbox{$\it t$}.  The kind of the new type constructor \mbox{$\it T$} is of the form
$\kappa_1\rightarrow\ldots\rightarrow\kappa_k\rightarrow\kappa$ where
the kinds \mbox{$\it \kappa_i$} of the arguments \mbox{$\it u_i$} and \mbox{$\it \kappa$} of the right hand
side \mbox{$\it t$} are determined by kind inference as described in
Section~\ref{kindinference}\index{kind}\index{kind inference}.
For example, the following definition can be used to provide an alternative
way of writing the list type constructor: 
\bprog
{\small\begin{verbatim}
  type List = []
\end{verbatim}}
\eprog
Type constructor symbols \mbox{$\it T$} introduced by type synonym declarations cannot
be partially applied; it is a static error to use \mbox{$\it T$} without the full number
of arguments.

Although recursive and mutually recursive datatypes are allowed,
\index{recursive datatype}
\index{type synonym!recursive}
this is not so for type synonyms, {\em unless an algebraic datatype
intervenes}.  For example,
\bprog
{\small\begin{verbatim}
  type Rec a   =  [Circ a]
  data Circ a  =  Tag [Rec a]
\end{verbatim}}
\eprog
is allowed, whereas
\bprog
{\small\begin{verbatim}
  type Rec a   =  [Circ a]        -- invalid
  type Circ a  =  [Rec a]         -- invalid
\end{verbatim}}
\eprog
is not. Similarly, \mbox{\tt type\ Rec\ a\ =\ [Rec\ a]} is not allowed.

Type synonyms are a convenient, but strictly syntactic, mechanism to make type
signatures more readable.  A synonym and its definition are completely
interchangeable, except in the instance type of an \mbox{\tt instance} declaration (Section~\ref{instance-decls}).

\subsection{Datatype Renamings}
\index{newtype declaration@{\tt newtype} declaration}
\label{datatype-renaming}

\begin{flushleft}\it\begin{tabbing}
\hspace{0.6in}\=\hspace{3.1in}\=\kill
$\it topdecl$\>\makebox[3.5em]{$\rightarrow$}$\it \makebox{\tt newtype}\ [context\ \makebox{\tt =>}]\ simpletype\ \makebox{\tt =}\ newconstr\ [deriving]$\\ 
$\it newconstr$\>\makebox[3.5em]{$\rightarrow$}$\it con\ atype$\\ 
$\it $\>\makebox[3.5em]{$|$}$\it con\ \makebox{\tt {\char'173}}\ var\ \makebox{\tt ::}\ type\ \makebox{\tt {\char'175}}$\\ 
$\it simpletype$\>\makebox[3.5em]{$\rightarrow$}$\it tycon\ tyvar_1\ \ldots \ tyvar_k$\>\makebox[3em]{}$\it (k\geq 0)$
\end{tabbing}\end{flushleft}
\index{topdecl@{\em topdecl} (\mbox{\tt newtype})}%
\indexsyn{simpletype}%
\indexsyn{newconstr}%

\noindent \noindent A declaration of the form
\[
\mbox{$\it \makebox{\tt newtype}\ cx\ \makebox{\tt =>}\ T\ u_1\ \ldots \ u_k\ \makebox{\tt =}\ N\ t$}
\]
introduces a new type whose
representation is the same as an existing type.  The type \mbox{$\it \makebox{\tt (}T\ u_1\\
\it \ldots \ u_k\makebox{\tt )}$} renames the datatype \mbox{$\it t$}.
It differs from a type synonym in
that it creates a distinct type that must be explicitly coerced to or
from the original type.  Also, unlike type synonyms, \mbox{\tt newtype} may be
used to define recursive types.
The constructor \mbox{$\it N$} in an expression 
coerces a value from type \mbox{$\it t$} to type \mbox{$\it \makebox{\tt (}T\ u_1\ \ldots \ u_k\makebox{\tt )}$}.
Using \mbox{$\it N$} in a pattern coerces a value from type \mbox{$\it \makebox{\tt (}T\ u_1\ \ldots \ u_k\makebox{\tt )}$}
to type \mbox{$\it t$}.  These coercions may be implemented without
execution time overhead; \mbox{\tt newtype} does not change the underlying
representation of an object.

New instances (see Section \ref{instance-decls}) can be defined for a
type defined by \mbox{\tt newtype} but may not be defined for a type synonym.  A type
created by \mbox{\tt newtype} differs from an algebraic datatype in that the
representation of an
algebraic datatype has an extra level of indirection.  This difference
may make access to the representation less efficient.  The difference is
reflected in different rules for pattern matching (see
Section~\ref{pattern-matching}).  Unlike algebraic datatypes, the
newtype constructor \mbox{$\it N$} is {\em unlifted}, so that \mbox{$\it N\ \bot$}
is the same as \mbox{$\it \bot$}.

The following examples clarify the differences between \mbox{\tt data} (algebraic
datatypes), \mbox{\tt type} (type synonyms), and \mbox{\tt newtype} (renaming types.)
Given the declarations 
\bprog
{\small\begin{verbatim}
  data D1 = D1 Int
  data D2 = D2 !Int
  type S = Int
  newtype N = N Int
  d1 (D1 i) = 42
  d2 (D2 i) = 42
  s i = 42
  n (N i) = 42
\end{verbatim}}
\eprog
the expressions \mbox{$\it \makebox{\tt (d1}\ \bot\makebox{\tt )}$}, \mbox{$\it \makebox{\tt (d2}\ \bot\makebox{\tt )}$} and 
\mbox{$\it \makebox{\tt (d2\ (D2}\ \bot\makebox{\tt ))}$} are all
equivalent to \mbox{$\it \bot$}, whereas \mbox{$\it \makebox{\tt (n}\ \bot\makebox{\tt )}$}, \mbox{$\it \makebox{\tt (n\ (N}\ \bot\makebox{\tt ))}$},
\mbox{$\it \makebox{\tt (d1\ (D1}\ \bot\makebox{\tt ))}$} and \mbox{$\it \makebox{\tt (s}\ \bot\makebox{\tt )}$}
are all equivalent to \mbox{\tt 42}.  In particular, \mbox{$\it \makebox{\tt (N}\ \bot\makebox{\tt )}$} is equivalent to
\mbox{$\it \bot$} while \mbox{$\it \makebox{\tt (D1}\ \bot\makebox{\tt )}$} is not equivalent to \mbox{$\it \bot$}.

The optional deriving part of a \mbox{\tt newtype} declaration is treated in
the same way as the deriving component of a \mbox{\tt data} declaration; see
Section~\ref{derived-decls}.

A \mbox{\tt newtype} declaration may use field-naming syntax, though of course
there may only be one field.  Thus:
\bprog
{\small\begin{verbatim}
  newtype Age = Age { unAge :: Int }
\end{verbatim}}
\eprog
brings into scope both a constructor and a de-constructor:
\bprog
{\small\begin{verbatim}
  Age   :: Int -> Age
  unAge :: Age -> Int
\end{verbatim}}
\eprog

\section{Type Classes and Overloading}
\index{class}
\index{overloading}
\label{overloading}
\label{classes}

\subsection{Class Declarations}
\index{class declaration}
%\index{class declaration@{\tt class} declaration}
\label{class-decls}

\begin{flushleft}\it\begin{tabbing}
\hspace{0.6in}\=\hspace{3.1in}\=\kill
$\it topdecl$\>\makebox[3.5em]{$\rightarrow$}$\it \makebox{\tt class}\ [scontext\ \makebox{\tt =>}]\ tycls\ tyvar\ [\makebox{\tt where}\ cdecls]$\\ 
$\it scontext$\>\makebox[3.5em]{$\rightarrow$}$\it simpleclass$\\ 
$\it $\>\makebox[3.5em]{$|$}$\it \makebox{\tt (}\ simpleclass_1\ \makebox{\tt ,}\ \ldots \ \makebox{\tt ,}\ simpleclass_n\ \makebox{\tt )}$\>\makebox[3em]{}$\it (n\geq 0)$\\ 
$\it simpleclass$\>\makebox[3.5em]{$\rightarrow$}$\it qtycls\ tyvar$\\ 
$\it cdecls$\>\makebox[3.5em]{$\rightarrow$}$\it \makebox{\tt {\char'173}}\ cdecl_1\ \makebox{\tt ;}\ \ldots \ \makebox{\tt ;}\ cdecl_n\ \makebox{\tt {\char'175}}$\>\makebox[3em]{}$\it (n\geq 0)$\\ 
$\it cdecl$\>\makebox[3.5em]{$\rightarrow$}$\it gendecl$\\ 
$\it $\>\makebox[3.5em]{$|$}$\it (funlhs\ |\ var)\ rhs$
\end{tabbing}\end{flushleft}
\index{topdecl@{\em topdecl} (\mbox{\tt class})}%
\indexsyn{gendecl}%
\indexsyn{cdecls}%
\indexsyn{cdecl}%
\index{declaration!within a {\tt class} declaration}

\noindent
A {\em class declaration} introduces a new class and the operations
({\em class methods}) on it.
A class declaration has the general form:
\[\ba{rl}
\mbox{$\it \makebox{\tt class}\ cx\ \makebox{\tt =>}\ C\ u\ \makebox{\tt where}\ cdecls$}
% \mbox{$\it \makebox{\tt {\char'173}}$}&\mbox{$\it \ldots \makebox{\tt ;}\ v_i\ \makebox{\tt ::}\ cx_i\ \makebox{\tt =>}\ t_i\ \makebox{\tt ;}\ \ldots \makebox{\tt ;}$}\\
%                               &\mbox{$\it \ldots \makebox{\tt ;}\ \makebox{\tt infix}\ n\ v_i\makebox{\tt ;}\ \ldots \makebox{\tt ;}$}\\
%                               &\mbox{$\it valdef_1\ \makebox{\tt ;}\ \ldots \ \makebox{\tt ;}\ valdef_m\ \makebox{\tt {\char'175}}$}
\ea\]
This introduces a new class name \mbox{$\it C$}; the type variable \mbox{$\it u$} is
scoped only over the class method signatures in the class body.
The context \mbox{$\it cx$} specifies the superclasses\index{superclass} of \mbox{$\it C$}, as
described below; the only type variable that may be referred to in \mbox{$\it cx$}
is \mbox{$\it u$}.

The superclass relation must not be cyclic; i.e.~it must form a
directed acyclic graph.\index{superclass}

The \mbox{$\it cdecls$} part of a \mbox{\tt class} declaration contains three kinds
of declarations:
\begin{itemize}
\item
The class declaration introduces new {\em class methods}
\index{class method}
\mbox{$\it v_i$}, whose scope extends outside the \mbox{\tt class} declaration.
The class methods of a class declaration are precisely the \mbox{$\it v_i$} for
which there is an explicit type signature
\[
\mbox{$\it v_i\ \makebox{\tt ::}\ cx_i\ \makebox{\tt =>}\ t_i$}
\]
in \mbox{$\it cdecls$}.  
Class methods share the top level namespace with variable
bindings and field names; they must not conflict with other top level
bindings in scope. 
That is, a class method can 
not have the same name as a top level definition, a field name, or
another class method.

The type of the top-level class method \mbox{$\it v_i$} is:
\[
v_i :: \forall u,\overline{w}.~(C u, cx_i) \Rightarrow t_i
\]
The \mbox{$\it t_i$} must mention \mbox{$\it u$}; it may mention type variables
\mbox{$\it \overline{w}$} other than \mbox{$\it u$}, in which case the type of \mbox{$\it v_i$} is
polymorphic in both \mbox{$\it u$} and \mbox{$\it \overline{w}\/$}.
The \mbox{$\it cx_i$} may constrain only \mbox{$\it \overline{w}$}; in particular,
the \mbox{$\it cx_i$} may not constrain \mbox{$\it u$}.
For example:
\bprog
{\small\begin{verbatim}
  class Foo a where
    op :: Num b => a -> b -> a
\end{verbatim}}
\eprog
Here the type of \mbox{\tt op} is
\mbox{$\it \forall\ a,\ b.~(\makebox{\tt Foo}~a,~\makebox{\tt Num}~b)~\ \Rightarrow\ a\ \rightarrow\ b\ \rightarrow\ a$}.

\item
The \mbox{$\it cdecls$} may also contain a {\em fixity declaration} for any of the class methods 
(but for no other values).
\index{fixity declaration}
However, since class methods declare top-level values, the fixity declaration for a class
method may alternatively appear at top level, outside the class declaration.

\item
Lastly, the \mbox{$\it cdecls$} may contain a
{\em default class method}
\index{default class method}
for any of the \mbox{$\it v_i$}.  The default class method for \mbox{$\it v_i$} is used if no binding for it
is given in a particular \mbox{\tt instance} declaration (see
Section~\ref{instance-decls}).
The default method declaration is a normal value definition, except that the
left hand side may only be a variable or function definition.  For example:
\bprog
{\small\begin{verbatim}
  class Foo a where
    op1, op2 :: a -> a
    (op1, op2) = ...
\end{verbatim}}
\eprog
is not permitted, because the left hand side of the default declaration is a
pattern.
\end{itemize}
Other than these cases, no other declarations are permitted in \mbox{$\it cdecls$}.

%Figure~\ref{standard-classes} shows some standard \Haskell{} classes,
%including the use of superclasses.
%; note the class inclusion diagram on the right.
%For example, \mbox{\tt Eq} is a superclass of \mbox{\tt Ord}, and thus in
%any context \mbox{\tt Ord\ a} is equivalent to \mbox{\tt (Eq\ a,\ Ord\ a)}.  Figures~\ref{data-class},
%\ref{prelude-bool} and \ref{numeric-inclusions} show the remaining
%standard \Haskell{} classes.

A \mbox{\tt class}
declaration with no \mbox{\tt where} part
\index{class declaration!with an empty \mbox{\tt where} part}
may be useful for combining a
collection of classes into a larger one that inherits all of the
class methods in the original ones.  For example:
\bprog
{\small\begin{verbatim}
  class  (Read a, Show a) => Textual a
\end{verbatim}}
\eprog
In such a case, if a type is an instance of all
superclasses,\index{superclass} it is 
not {\em automatically} an instance of the subclass, even though the
subclass has no immediate class methods.  The \mbox{\tt instance} declaration must be
given explicitly with no \mbox{\tt where} part.
\index{instance declaration!with an empty \mbox{\tt where} part}

\subsection{Instance Declarations}
\label{instance-decls}
\index{instance declaration}

\begin{flushleft}\it\begin{tabbing}
\hspace{0.6in}\=\hspace{3.1in}\=\kill
$\it topdecl$\>\makebox[3.5em]{$\rightarrow$}$\it \makebox{\tt instance}\ [scontext\ \makebox{\tt =>}]\ qtycls\ inst\ [\makebox{\tt where}\ idecls]$\\ 
$\it inst$\>\makebox[3.5em]{$\rightarrow$}$\it gtycon$\\ 
$\it $\>\makebox[3.5em]{$|$}$\it \makebox{\tt (}\ gtycon\ tyvar_1\ \ldots \ tyvar_k\ \makebox{\tt )}$\>\makebox[3em]{}$\it (k\geq 0,\ tyvars\ {\rm\ distinct})$\\ 
$\it $\>\makebox[3.5em]{$|$}$\it \makebox{\tt (}\ tyvar_1\ \makebox{\tt ,}\ \ldots \ \makebox{\tt ,}\ tyvar_k\ \makebox{\tt )}$\>\makebox[3em]{}$\it (k\geq 2,\ tyvars\ {\rm\ distinct})$\\ 
$\it $\>\makebox[3.5em]{$|$}$\it \makebox{\tt [}\ tyvar\ \makebox{\tt ]}$\\ 
$\it $\>\makebox[3.5em]{$|$}$\it \makebox{\tt (}\ tyvar_1\ \makebox{\tt ->}\ tyvar_2\ \makebox{\tt )}$\>\makebox[3em]{}$\it (tyvar_1\ {\rm\ and}\ tyvar_2\ {\rm\ distinct})$\\ 
$\it idecls$\>\makebox[3.5em]{$\rightarrow$}$\it \makebox{\tt {\char'173}}\ idecl_1\ \makebox{\tt ;}\ \ldots \ \makebox{\tt ;}\ idecl_n\ \makebox{\tt {\char'175}}$\>\makebox[3em]{}$\it (n\geq 0)$\\ 
$\it idecl$\>\makebox[3.5em]{$\rightarrow$}$\it (funlhs\ |\ var)\ rhs$\\ 
$\it $\>\makebox[3.5em]{$|$}$\it $\>\makebox[3em]{}$\it (empty)$
\end{tabbing}\end{flushleft}
\index{topdecl@{\em topdecl} (\mbox{\tt instance})}
\indexsyn{inst}%
\indexsyn{valdefs}%
\indexsyn{gtycon}%
\indexsyn{idecl}%
\indexsyn{idecls}%
%\index{instance declaration@{\tt instance} declaration}
\index{declaration!within an {\tt instance} declaration}
%\ToDo{tycls left off above}
An {\em instance declaration} introduces an instance of a class.  Let
\[ \mbox{$\it \makebox{\tt class}\ cx\ \makebox{\tt =>}\ C\ u\ \makebox{\tt where}\ \makebox{\tt {\char'173}}\ cbody\ \makebox{\tt {\char'175}}$} \]
be a \mbox{\tt class} declaration.  The general form of the corresponding
instance declaration is:
\[ \mbox{$\it \makebox{\tt instance}\ cx'\ \makebox{\tt =>}\ C\ (T\ u_1\ \ldots \ u_k)\ \makebox{\tt where}\ \makebox{\tt {\char'173}}\ d\ \makebox{\tt {\char'175}}$} \]
where \mbox{$\it k\geq0$}.
The type \mbox{$\it (T\ u_1\ \ldots \ u_k)$} must take the form of
a type constructor \mbox{$\it T$} applied to simple type variables \mbox{$\it u_1,\ \ldots \ u_k$};
furthermore, \mbox{$\it T$} must not be a type synonym\index{type synonym}, 
and the \mbox{$\it u_i$} must all be distinct.

This prohibits instance declarations
such as:
\bprog
{\small\begin{verbatim}
  instance C (a,a) where ...
  instance C (Int,a) where ...
  instance C [[a]] where ...
\end{verbatim}}
\eprog
The declarations \mbox{$\it d$} may contain bindings only for the class
methods\index{class method} of \mbox{$\it C$}.  It is illegal to give a 
binding for a class method that is not in scope, but the name under
which it is in scope is immaterial; in particular, it may be a qualified
name.  (This rule is identical to that used for subordinate names in
export lists --- Section~\ref{export}.)
For example, this is legal, even though \mbox{\tt range} is in scope only
with the qualified name \mbox{\tt Data.Ix.range}.
\bprog
{\small\begin{verbatim}
  module A where
    import qualified Data.Ix

    instance Data.Ix.Ix T where
      range = ...
\end{verbatim}}
\eprog
The declarations may not contain any type
signatures or fixity declarations,\index{type signature}\index{fixity declaration}
since these have already been given in the \mbox{\tt class}
declaration.  As in the case of default class methods
\index{default class method}%
(Section~\ref{class-decls}), the method declarations must take the form of
a variable or function definition.

% However, unlike other declarations, the name of the bound
% variable may be qualified.  So this is legal:
% \bprog
% 
%  module A where
%    import qualified B( Foo(op) )
%    instance B.Foo Int where
%      B.op = ...
%
%\eprog
% Here, module \mbox{\tt A} imports class \mbox{\tt Foo} from module \mbox{\tt B}, and then gives an 
%instance declaration for \mbox{\tt Foo}.  Since \mbox{\tt B} is imported \mbox{\tt qualified}, the name
%of the class method, \mbox{\tt op}, is not in scope; so the definition must define \mbox{\tt B.op}.
% Hence the need for the \mbox{$\it qfunlhs$} and \mbox{$\it qvar$} left hand sides in \mbox{$\it idecl$}.

If no binding is given for some class method then the
corresponding default class method
\index{default class method}
in the \mbox{\tt class} declaration is used (if
present); if such a default does
not exist then the class method of this instance
is bound to \mbox{\tt undefined} 
and no compile-time error results.

An \mbox{\tt instance} declaration that makes the type \mbox{$\it T$} to be an instance
of class \mbox{$\it C$} is called a {\em C-T instance declaration}
\index{instance declaration} and is
subject to these static restrictions:
%\index{instance declaration!with respect to modules}
\begin{itemize}

\item A type may not be declared as an instance of a
particular class more than once in the program.

\item The class and type must have the same kind; 
this can be determined using kind inference as described
in Section~\ref{kindinference}\index{kind}\index{kind inference}.

\item 
Assume that the type variables in the instance type \mbox{$\it (T\ u_1\ \ldots \ u_k)$}
satisfy the constraints in the instance context \mbox{$\it cx'$}.  Under this
assumption, the following two conditions must also be satisfied:
\begin{enumerate}
\item
The constraints expressed by the superclass context \mbox{$\it cx[(T\ u1\ \ldots \ uk)/u]$}
        of \mbox{$\it C$} must be satisfied.  In other words, \mbox{$\it T$} must be an instance
        of each of \mbox{$\it C$}'s superclasses and the contexts of all
        superclass instances must be implied by \mbox{$\it cx'$}.
\item
Any constraints on the type variables in the instance type
      that are required for the class method declarations in \mbox{$\it d$} to be
      well-typed must also be satisfied.
\end{enumerate}

% \mbox{$\it T$} must be an instance of each of \mbox{$\it C$}'s superclasses.  More precisely,
% under the assumption that the constraints expressed by the
% instance context \mbox{$\it cx'$} are satisfied, the constraints expressed by the
% superclass context \mbox{$\it cx[(T\ u1\ \ldots \ uk)/u]$} must also be satisfied.
% 
% Furthermore, under the same assumption, any constraints arising from the
% class method declarations in \mbox{$\it d$} must also be satisfied.  

In fact, except in pathological cases 
it is possible to infer from the instance declaration the
most general instance context \mbox{$\it cx'$} satisfying the above two constraints, 
but it is nevertheless mandatory
to write an explicit instance context.
\end{itemize}
The following example illustrates the restrictions imposed by superclass instances:
\bprog
{\small\begin{verbatim}
  class Foo a => Bar a where ...
  
  instance (Eq a, Show a) => Foo [a] where ...
  
  instance Num a => Bar [a] where ...
\end{verbatim}}
\eprog
This example is valid Haskell.  Since \mbox{\tt Foo} is a superclass of \mbox{\tt Bar},
the second instance declaration is only valid if \mbox{\tt [a]} is an
instance of \mbox{\tt Foo} under the assumption \mbox{\tt Num\ a}.  
The first instance declaration does indeed say that \mbox{\tt [a]} is an instance
of \mbox{\tt Foo} under this assumption, because \mbox{\tt Eq} and \mbox{\tt Show} are superclasses
of \mbox{\tt Num}.

If the two instance declarations instead read like this:
\bprog
{\small\begin{verbatim}
  instance Num a => Foo [a] where ...
  
  instance (Eq a, Show a) => Bar [a] where ...
\end{verbatim}}
\eprog
then the program would be invalid.  The second instance declaration is
valid only if \mbox{\tt [a]} is an instance of \mbox{\tt Foo} under the assumptions
\mbox{\tt (Eq\ a,\ Show\ a)}.  But this does not hold, since \mbox{\tt [a]} is only an
instance of \mbox{\tt Foo} under the stronger assumption \mbox{\tt Num\ a}.

Further examples of 
\mbox{\tt instance} declarations may be found in Chapter~\ref{stdprelude}.

\subsection{Derived Instances}
\index{derived instance}
\label{derived-decls}

As mentioned in Section~\ref{datatype-decls}, \mbox{\tt data} and \mbox{\tt newtype}
declarations 
contain an optional \mbox{\tt deriving} form.  If the form is included, then
{\em derived instance declarations} are automatically generated for
the datatype in each of the named classes.
These instances are subject to the same restrictions as user-defined
instances.  When deriving a class \mbox{$\it C$} for a type \mbox{$\it T$}, instances for
all superclasses of \mbox{$\it C$} must exist for \mbox{$\it T$}, either via an explicit
\mbox{\tt instance} declaration or by including the superclass in the
\mbox{\tt deriving} clause.

Derived instances provide convenient commonly-used operations for
user-de\-fined da\-ta\-types.  For example, derived instances for datatypes
in the class \mbox{\tt Eq} define the operations \mbox{\tt ==} and \mbox{\tt /=}, freeing the
programmer from the need to define them.
%and taking advantage of
%\Haskell{}'s class mechanism to overload these operations.

The only classes in the Prelude for
which derived instances are allowed are
\mbox{\tt Eq}\indexdi{Eq},
\mbox{\tt Ord}\indexdi{Ord},
\mbox{\tt Enum}\indexdi{Enum},
\mbox{\tt Bounded}\indexdi{Bounded},
\mbox{\tt Show}\indexdi{Show},
and \mbox{\tt Read}\indexdi{Read},
all mentioned in Figure~\ref{standard-classes}.
The
precise details of how the derived instances are generated for each of
these classes are provided in Chapter~\ref{derived-appendix}, including
a specification of when such derived instances are possible. 
%(which is important for the following discussion).
Classes defined by the standard libraries may also be derivable.

A static error results if it is not possible to derive an \mbox{\tt instance}
declaration over a class named in a \mbox{\tt deriving} form.  For example,
not all datatypes can properly support class methods in
\mbox{\tt Enum}.\indexclass{Enum}  It is 
also a static error to give an explicit \mbox{\tt instance} declaration for
a class that is also derived.
%These rules also apply to the superclasses
%of the class in question.

If the \mbox{\tt deriving} form is omitted from a \mbox{\tt data} or \mbox{\tt newtype}
declaration, then {\em no} instance declarations
are derived for 
that datatype; that is, omitting a \mbox{\tt deriving} form is equivalent to
including an empty deriving form: \mbox{\tt deriving\ ()}.

% OLD:
%On the other hand, if the \mbox{\tt deriving} form is omitted from a \mbox{\tt data}
%declaration, then \mbox{\tt instance} declarations are derived for the datatype
%in as many of the six classes mentioned above as is possible (see
%Chapter~\ref{derived-appendix}); that is, no
%static error will result if the \mbox{\tt instance} declarations cannot be generated.

%OLD:
%If {\em no} derived \mbox{\tt instance} declarations for a datatype
%are wanted, then the empty deriving form \mbox{\tt deriving\ ()} must be given
%in the \mbox{\tt data} declaration for that type.

\subsection{Ambiguous Types, and Defaults for Overloaded Numeric Operations}
\label{default-decls}
\index{default declaration@{\tt default} declaration}
\index{overloading!defaults}

\begin{flushleft}\it\begin{tabbing}
\hspace{0.6in}\=\hspace{3.1in}\=\kill
$\it topdecl$\>\makebox[3.5em]{$\rightarrow$}$\it \makebox{\tt default}\ \makebox{\tt (}type_1\ \makebox{\tt ,}\ \ldots \ \makebox{\tt ,}\ type_n\makebox{\tt )}$\>\makebox[3em]{}$\it (n\geq 0)$
\end{tabbing}\end{flushleft}
\index{topdecl@{\em topdecl} (\mbox{\tt default})}

\noindent
A problem inherent with \Haskell{}-style overloading is the
possibility of an {\em ambiguous type}.
\index{ambiguous type}
For example, using the
\mbox{\tt read} and \mbox{\tt show} functions defined in Chapter~\ref{derived-appendix},
and supposing that just \mbox{\tt Int} and \mbox{\tt Bool} are members of \mbox{\tt Read} and
\mbox{\tt Show}, then the expression
\bprog
{\small\begin{verbatim}
  let x = read "..." in show x  -- invalid
\end{verbatim}}
\eprog
is ambiguous, because the types for \mbox{\tt show} and \mbox{\tt read},
\[\ba{ll}
\mbox{\tt show} & \mbox{$\it ::\ \forall\ a.~\makebox{\tt Show}~\ a\ \Rightarrow\ a\ \rightarrow\ \makebox{\tt String}$} \\
\mbox{\tt read} & \mbox{$\it ::\ \forall\ a.~\makebox{\tt Read}~\ a\ \Rightarrow\ \makebox{\tt String}\ \rightarrow\ a$}
\ea\]
could be satisfied by instantiating \mbox{\tt a} as either \mbox{\tt Int}
in both cases, or \mbox{\tt Bool}.  Such expressions
are considered ill-typed, a static error.

We say that an expression \mbox{\tt e} has an {\em ambiguous type}
if, in its type \mbox{$\it \forall\ \overline{u}.~cx\ \Rightarrow\ t$}, 
there is a type variable \mbox{$\it u$} in \mbox{$\it \overline{u}$} that occurs in \mbox{$\it cx$} 
but not in \mbox{$\it t$}.  Such types are invalid.

For example, the earlier expression involving \mbox{\tt show} and \mbox{\tt read} has
an ambiguous type since its type is 
\mbox{$\it \forall\ a.~\ \makebox{\tt Show}~\ a,\ \makebox{\tt Read}~\ a\ \Rightarrow\ \makebox{\tt String}$}.

Ambiguous types can only be circumvented by
input from the user.  One way is through the use of {\em expression
type-signatures}
\index{expression type-signature}
as described in Section~\ref{expression-type-sigs}.
For example, for the ambiguous expression given earlier, one could
write:
\bprog
{\small\begin{verbatim}
  let x = read "..." in show (x::Bool)
\end{verbatim}}
\eprog
which disambiguates the type.

Occasionally, an otherwise ambiguous expression needs to be made
the same type as some variable, rather than being given a fixed
type with an expression type-signature.  This is the purpose
of the function \mbox{\tt asTypeOf} (Chapter~\ref{stdprelude}):
\mbox{$\it x$} \mbox{\tt `asTypeOf`} \mbox{$\it y$} has the value of \mbox{$\it x$}, but \mbox{$\it x$} and \mbox{$\it y$} are
forced to have the same type.  For example,
\bprog
{\small\begin{verbatim}
  approxSqrt x = encodeFloat 1 (exponent x `div` 2) `asTypeOf` x
\end{verbatim}}
\eprog
(See Section~\ref{coercion} for a description of \mbox{\tt encodeFloat} and \mbox{\tt exponent}.)

Ambiguities in the class \mbox{\tt Num}\indexclass{Num}
are most common, so \Haskell{}
provides another way to resolve them---with a {\em
default declaration}:
\[
\mbox{$\it \makebox{\tt default\ (}t_1\ \makebox{\tt ,}\ \ldots \ \makebox{\tt ,}\ t_n\makebox{\tt )}$}
\]
where \mbox{$\it n\geq0$}, and each
\mbox{$\it t_i$} must be a type for which \mbox{$\it \makebox{\tt Num\ }t_i$} holds.
In situations where an ambiguous type is discovered, an ambiguous type variable, \mbox{$\it v$}, 
is defaultable if:
\begin{itemize}
\item \mbox{$\it v$} appears only in constraints of the form \mbox{$\it C\ v$}, where \mbox{$\it C$} is a class, and
\item at least one of these classes is a numeric class,
        (that is, \mbox{\tt Num} or a subclass of \mbox{\tt Num}), and 
\item all of these classes are defined in the Prelude or a standard library
(Figures~\ref{basic-numeric-1}--\ref{basic-numeric-2}
show the numeric classes, and
Figure~\ref{standard-classes}
shows the classes defined in the Prelude.)
\end{itemize}
Each defaultable variable is replaced by the first type in the
default list that is an instance of all the ambiguous variable's classes.
It is a static error if no such type is found.

Only one default declaration is permitted per module, and its effect
is limited to that module.  If no default declaration is given in a
module then it assumed to be:
\bprog
{\small\begin{verbatim}
  default (Integer, Double)
\end{verbatim}}
\eprog
The empty default declaration, \mbox{\tt default\ ()}, turns off all defaults in a module.

\section{Nested Declarations}
\label{nested}

The following declarations may be used in any declaration list,
including the top level of a module.

\subsection{Type Signatures}
\index{type signature}
\label{type-signatures}

\begin{flushleft}\it\begin{tabbing}
\hspace{0.6in}\=\hspace{3.1in}\=\kill
$\it gendecl$\>\makebox[3.5em]{$\rightarrow$}$\it vars\ \makebox{\tt ::}\ [context\ \makebox{\tt =>}]\ type$\\ 
$\it vars$\>\makebox[3.5em]{$\rightarrow$}$\it var_1\ \makebox{\tt ,}\ \ldots \makebox{\tt ,}\ var_n$\>\makebox[3em]{}$\it (n\geq 1)$
\end{tabbing}\end{flushleft}
\indexsyn{signdecl}%
\indexsyn{vars}%
A type signature specifies types for variables, possibly with respect
to a context.  A type signature has the form:
\[
\mbox{$\it v_1,\ \ldots ,\ v_n\ \makebox{\tt ::}\ cx\ \makebox{\tt =>}\ t$}
\]
which is equivalent to asserting
\mbox{$\it v_i\ \makebox{\tt ::}\ cx\ \makebox{\tt =>}\ t$}
for each \mbox{$\it i$} from \mbox{$\it 1$} to \mbox{$\it n$}.  Each \mbox{$\it v_i$} must have a value binding in
the same declaration list that contains the type signature; i.e.~it is
invalid to give a type signature for a variable bound in an
outer scope.
Moreover, it is invalid to give more than one type signature for one
variable, even if the signatures are identical.

As mentioned in Section~\ref{type-syntax},
every type variable appearing in a signature
is universally quantified over that signature, and hence
the scope of a type variable is limited to the type
signature that contains it.  For example, in the following
declarations
\bprog
{\small\begin{verbatim}
  f :: a -> a
  f x = x :: a                  -- invalid
\end{verbatim}}
\eprog
the \mbox{\tt a}'s in the two type signatures are quite distinct.  Indeed,
these declarations contain a static error, since \mbox{\tt x} does not have
type \mbox{$\it \forall\ a.~a$}.  (The type of \mbox{\tt x} is dependent on the type of
\mbox{\tt f}; there is currently no way in \Haskell{} to specify a signature
for a variable with a dependent type; this is explained in Section
\ref{monomorphism}.)

If a given program includes a signature
for a variable \mbox{$\it f$}, then each use of \mbox{$\it f$} is treated as having the
declared type.  It is a static error if the same type cannot also be
inferred for the defining occurrence of \mbox{$\it f$}.

If a variable \mbox{$\it f$} is defined without providing a corresponding type
signature declaration, then each use of \mbox{$\it f$} outside its own declaration
group (see Section~\ref{dependencyanalysis}) is treated as having the
corresponding inferred, or {\em principal} type \index{principal type}.
However, to ensure that type inference is still possible, the defining
occurrence, and all uses of \mbox{$\it f$} within its declaration group must have
the same monomorphic type (from which the principal type is obtained
by generalization, as described in Section~\ref{generalization}).

%%A type signature for \mbox{$\it x$} may be more specific than the principal
%%type derivable from the value binding of \mbox{$\it x$} (see
%%Section~\ref{type-semantics}), but it is an error to give a type
%%that is more
%%general than, or incomparable to, the principal type.
%%\index{principal type}
%%If a more specific type is given then all occurrences of the
%%variable must be used at the more specific type or at a more
%%specific type still.
%
For example, if we define\nopagebreak[4]
\bprog
{\small\begin{verbatim}
  sqr x  =  x*x
\end{verbatim}}
\eprog
then the principal type is 
\mbox{$\it \makebox{\tt sqr}\ ::\ \forall\ a.~\ \makebox{\tt Num}~\ a\ \Rightarrow\ a\ \rightarrow\ a$}, 
which allows
applications such as \mbox{\tt sqr\ 5} or \mbox{\tt sqr\ 0.1}.  It is also valid to declare
a more specific type, such as
\bprog
{\small\begin{verbatim}
  sqr :: Int -> Int
\end{verbatim}}
\eprog
but now applications such as \mbox{\tt sqr\ 0.1} are invalid.  Type signatures such as
\bprog
{\small\begin{verbatim}
  sqr :: (Num a, Num b) => a -> b     -- invalid
  sqr :: a -> a                       -- invalid
\end{verbatim}}
\eprog
are invalid, as they are more general than the principal type of \mbox{\tt sqr}.

Type signatures can also be used to support
{\em polymorphic recursion}\index{polymorphic recursion}.
The following definition is pathological, but illustrates how a type
signature can be used to specify a type more general than the one that
would be inferred:
\bprog
{\small\begin{verbatim}
  data T a  =  K (T Int) (T a)
  f         :: T a -> a
  f (K x y) =  if f x == 1 then f y else undefined
\end{verbatim}}
\eprog
If we remove the signature declaration, the type of \mbox{\tt f} will be
inferred as \mbox{\tt T\ Int\ ->\ Int} due to the first recursive call for which
the argument to \mbox{\tt f} is \mbox{\tt T\ Int}.  Polymorphic recursion allows the user
to supply the more general type signature, \mbox{\tt T\ a\ ->\ a}.

\subsection{Fixity Declarations}
\index{fixity declaration}
\label{fixity}

\begin{flushleft}\it\begin{tabbing}
\hspace{0.6in}\=\hspace{3.1in}\=\kill
$\it gendecl$\>\makebox[3.5em]{$\rightarrow$}$\it fixity\ [integer]\ ops$\\ 
$\it fixity$\>\makebox[3.5em]{$\rightarrow$}$\it \makebox{\tt infixl}\ |\ \makebox{\tt infixr}\ |\ \makebox{\tt infix}$\\ 
$\it ops$\>\makebox[3.5em]{$\rightarrow$}$\it op_1\ \makebox{\tt ,}\ \ldots \ \makebox{\tt ,}\ op_n$\>\makebox[3em]{}$\it (n\geq 1)$\\ 
$\it op$\>\makebox[3.5em]{$\rightarrow$}$\it varop\ |\ conop$
\end{tabbing}\end{flushleft}
\indexsyn{gendecl}%
\indexsyn{fixity}%
\indexsyn{ops}%
\indexsyn{op}%
%\indextt{infixl}\indextt{infixr}\indextt{infix}
A fixity declaration gives the fixity and binding
precedence of one or more operators.  The \mbox{$\it integer$} in a fixity declaration
must be in the range \mbox{$\it 0$} to \mbox{$\it 9$}.
A fixity declaration may appear anywhere that 
a type signature appears and, like a type signature, declares a property of
a particular operator.  Also like a type signature,
a fixity declaration can only occur in the same sequence of declarations as
the declaration of the operator itself, and at most one fixity declaration
may be given for any operator.  (Class methods are a minor exception;
their fixity declarations can occur either in the class declaration itself
or at top level.)

There are three kinds of fixity, non-, left- and right-associativity
(\mbox{\tt infix}, \mbox{\tt infixl}, and \mbox{\tt infixr}, respectively), and ten precedence
levels, 0 to 9 inclusive (level 0 binds least tightly, and level 9
binds most tightly).  If the \mbox{$\it digit$} is omitted, level 9 is assumed.
Any operator lacking a fixity declaration
is assumed to be \mbox{\tt infixl\ 9} (See Section~\ref{expressions} for more on
the use of fixities).
Table~\ref{prelude-fixities} lists the fixities and precedences of
the operators defined in the Prelude.

\begin{table}[htb]
\[
\centerline{
\bto{|r||l|l|l|}
\hline 
Prec-  & Left associative       & Non-associative       & Right associative \\
edence & operators              & operators             & operators \\ \hline\hline
%9  & \mbox{\tt !}, \mbox{\tt !!}, \mbox{\tt //}   &                               & \mbox{\tt .}             \\ \hline
9  & \mbox{\tt !!}               &                               & \mbox{\tt .}             \\ \hline
8  &                    &                               & \mbox{\tt {\char'136}}, \mbox{\tt {\char'136}{\char'136}}, \mbox{\tt **} \\ \hline
7  & \mbox{\tt *}, \mbox{\tt /}, \mbox{\tt `div`},                 &               &  \\
   & \mbox{\tt `mod`}, \mbox{\tt `rem`}, \mbox{\tt `quot`}         &               &  \\ \hline
6  & \mbox{\tt +}, \mbox{\tt -}           &                               &                 \\ \hline
5  &                    &                               & \mbox{\tt :}, \mbox{\tt ++}       \\ \hline
4  &                    & \mbox{\tt ==}, \mbox{\tt /=}, \mbox{\tt <}, \mbox{\tt <=}, \mbox{\tt >}, \mbox{\tt >=}, &             \\
   &                    & \mbox{\tt `elem`}, \mbox{\tt `notElem`}         &                 \\ \hline
3  &                    &                               & \mbox{\tt {\char'46}{\char'46}}            \\ \hline
2  &                    &                               & \mbox{\tt ||}            \\ \hline
1  &    \mbox{\tt >>}, \mbox{\tt >>=}     &                               &               \\ \hline
0  &                    &                               & \mbox{\tt {\char'44}}, \mbox{\tt {\char'44}!}, \mbox{\tt `seq`}    \\ \hline
\eto}
\]
%**<div align=center> <h4>Table 2</h4> </div>
\ecaption{Precedences and fixities of prelude operators}
% Removed (:%) -- infixl 7, on the grounds that it's not exported from 
% PreludeRatio!
\label{prelude-fixities}
\index{"!"!@{\tt  {\char'041}{\char'041}}}
\indextt{.}
\indextt{**}
\index{^@{\tt  {\char'136}}} % this is ^
\index{^^@{\tt  {\char'136}{\char'136}}} % this is ^^
\indextt{*}
\indextt{/}
\indextt{div}
\indextt{mod}
\indextt{rem}
\indextt{+}
\indextt{-}
\indextt{:}
\indextt{++}
\indextt{/=}
\indextt{<}
\indextt{<=}
\indextt{==}
\indextt{>}
\indextt{>=}
\indextt{elem}
\indextt{notElem}
\index{&&@{\tt  \&\&}}
\index{\vert\vert@{\tt  {\char'174}{\char'174}}}
\indextt{>>}
\indextt{>>=}
\index{$@{\tt  {\char'044}}}
\end{table}

% Old table
% \begin{tabular}{|c|l|l|}\hline
% Precedence & Fixity & Operators \\ \hline
% 9 & \mbox{\tt infixl} & \mbox{\tt !}, \mbox{\tt !!}, \mbox{\tt //} \\
%   & \mbox{\tt infixr} & \mbox{\tt .} \\
% 8 & \mbox{\tt infixr} & \mbox{\tt **}, \mbox{\tt {\char'136}}, \mbox{\tt {\char'136}{\char'136}} \\
% 7 & \mbox{\tt infixl} & \mbox{\tt {\char'45}}, \mbox{\tt *} \\
%   & \mbox{\tt infix}  & \mbox{\tt /}, \mbox{\tt `div`}, \mbox{\tt `mod`}, \mbox{\tt `rem`} \\
% 6 & \mbox{\tt infixl} & \mbox{\tt +}, \mbox{\tt -} \\
%   & \mbox{\tt infix}  & \mbox{\tt :+} \\
% 5 & \mbox{\tt infixr} & \mbox{\tt :}, \mbox{\tt ++} \\
%   & \mbox{\tt infix}  & \mbox{\tt {\char'134}{\char'134}} \\
% 4 & \mbox{\tt infix}  & \mbox{\tt /=}, \mbox{\tt <}, \mbox{\tt <=}, \mbox{\tt ==}, \mbox{\tt >}, \mbox{\tt >=}, \mbox{\tt `elem`}, \mbox{\tt `notElem`} \\
% 3 & \mbox{\tt infixr} & \mbox{\tt {\char'46}{\char'46}} \\
% 2 & \mbox{\tt infixr} & \mbox{\tt ||} \\
% 1 & \mbox{\tt infixr}  & \mbox{\tt >>}, \mbox{\tt >>=} \\
% 0 & \mbox{\tt infixr} & \mbox{\tt {\char'44}} \\ \hline
% \end{tabular}}

Fixity is a property of a particular entity (constructor or variable), just like
its type; fixity is not a property of that entity's {\em name}.
For example: 
\bprog
{\small\begin{verbatim}
  module Bar( op ) where
    infixr 7 `op`
    op = ...
  
  module Foo where
    import qualified Bar
    infix 3 `op`
  
    a `op` b = (a `Bar.op` b) + 1
  
    f x = let
             p `op` q = (p `Foo.op` q) * 2
          in ...
\end{verbatim}}
\eprog
Here, \mbox{\tt `Bar.op`} is \mbox{\tt infixr\ 7}, \mbox{\tt `Foo.op`} is \mbox{\tt infix\ 3}, and
the nested definition of \mbox{\tt op} in \mbox{\tt f}'s right-hand side has the
default fixity of \mbox{\tt infixl\ 9}.  (It would also be possible
to give a fixity to the nested definition of \mbox{\tt `op`} with a nested
fixity declaration.)



\subsection{Function and Pattern Bindings}
 \label{function-bindings}\label{pattern-bindings}
\index{function binding}\index{pattern binding}

\begin{flushleft}\it\begin{tabbing}
\hspace{0.6in}\=\hspace{3.1in}\=\kill
$\it decl$\>\makebox[3.5em]{$\rightarrow$}$\it (funlhs\ |\ \hprime{pat})\ rhs$\\ 
$\it $\\ 
$\it funlhs$\>\makebox[3.5em]{$\rightarrow$}$\it var\ apat\ \{\ apat\ \}$\\ 
$\it $\>\makebox[3.5em]{$|$}$\it \hprime{pat\ varop\ pat}$\\ 
$\it $\>\makebox[3.5em]{$|$}$\it \makebox{\tt (}\ funlhs\ \makebox{\tt )}\ apat\ \{\ apat\ \}$\\ 
$\it $\\ 
$\it rhs$\>\makebox[3.5em]{$\rightarrow$}$\it \makebox{\tt =}\ exp\ [\makebox{\tt where}\ decls]$\\ 
$\it $\>\makebox[3.5em]{$|$}$\it gdrhs\ [\makebox{\tt where}\ decls]$\\ 
$\it $\\ 
$\it gdrhs$\>\makebox[3.5em]{$\rightarrow$}$\it \hprime{guards\ \makebox{\tt =}\ exp\ [gdrhs]}$\\ 
$\it $\\ 
$\it \hprime{guards}$\>\makebox[3.5em]{$\rightarrow$}$\it \hprime{\makebox{\tt |}\ guard_1,\ \ldots ,\ guard_n}$\>\makebox[3em]{}$\it \hprime{(n\geq 1)}$\\ 
$\it $\\ 
$\it \indexsyn{guard}\hprime{guard}$\>\makebox[3.5em]{$\rightarrow$}$\it \hprime{pat\ \makebox{\tt <-}\ infixexp}$\>\makebox[3em]{}$\it (\hprime{\tr{pattern\ guard}})$\\ 
$\it $\>\makebox[3.5em]{$|$}$\it \hprime{\makebox{\tt let}\ decls}$\>\makebox[3em]{}$\it (\hprime{\tr{local\ declaration}})$\\ 
$\it $\>\makebox[3.5em]{$|$}$\it infixexp$\>\makebox[3em]{}$\it (\tr{boolean\ guard})$
\end{tabbing}\end{flushleft}
\indexsyn{decl}%
\indexsyn{pat}%
\indexsyn{rhs}%
\indexsyn{gdrhs}%
\indexsyn{guards}%
\indexsyn{guard}%

We distinguish two cases within this syntax: a {\em pattern binding}
occurs when the left hand side is a \hprime{\mbox{$\it pat$}}; 
otherwise, the binding is called a {\em function
binding}.  Either binding may appear at the top-level of a module or
within a \mbox{\tt where} or \mbox{\tt let} construct.  

\subsubsection{Function bindings}
\index{function binding}
A function binding binds a variable to a function value.  The general
form of a function binding for variable \mbox{$\it x$} is:
\[\ba{lll}
\mbox{$\it x$} & \mbox{$\it p_{11}\ \ldots \ p_{1k}$} & \mbox{$\it match_1$}\\
\mbox{$\it \ldots $} \\
\mbox{$\it x$} & \mbox{$\it p_{n1}\ \ldots \ p_{nk}$} & \mbox{$\it match_n$}
\ea\]
where each \mbox{$\it p_{ij}$} is a pattern, and where each \mbox{$\it match_i$} is of the
general form:
\[\ba{l}
\mbox{$\it \makebox{\tt =}\ e_i\ \makebox{\tt where\ {\char'173}}\ decls_i\ \makebox{\tt {\char'175}}$}
\ea\]
or
\[\ba{lll}

\mbox{$\it \makebox{\tt |}\ \hprime{gs_{i1}}$}   & \mbox{$\it \makebox{\tt =}\ e_{i1}$} \\
\mbox{$\it \ldots $} \\
\mbox{$\it \makebox{\tt |}\ \hprime{gs_{im_i}}$} & \mbox{$\it \makebox{\tt =}\ e_{im_i}$} \\
               & \multicolumn{2}{l}{\mbox{$\it \makebox{\tt where\ {\char'173}}\ decls_i\ \makebox{\tt {\char'175}}$}}
\ea\]
and where \mbox{$\it n\geq 1$}, \mbox{$\it 1\leq i\leq n$}, \mbox{$\it m_i\geq 1$}.  The former is treated
as shorthand for a particular case of the latter, namely:
\[\ba{l}
\mbox{$\it \makebox{\tt |\ True\ =}\ e_i\ \makebox{\tt where\ {\char'173}}\ decls_i\ \makebox{\tt {\char'175}}$}
\ea\]

Note that all clauses defining a function must be contiguous, and the
number of patterns in each clause must be the same.  The set of
patterns corresponding to each match must be {\em
linear}\index{linearity}\index{linear pattern}---no variable is
allowed to appear more than once in the entire set.

Alternative syntax is provided for binding functional values to infix
operators.  For example, these three function
definitions are all equivalent:
\bprog
{\small\begin{verbatim}
  plus x y z = x+y+z
  x `plus` y = \ z -> x+y+z
  (x `plus` y) z = x+y+z
\end{verbatim}}
\eprog

\begin{haskellprime}

Note that fixity resolution applies to the infix variants of the
function binding in the same way as for expressions
(Section~\ref{fixity-resolution}).  Applying fixity resolution to the
left side of the equals in a function binding must leave the \mbox{$\it varop$}
being defined at the top level.  For example, if we are defining a new
operator \mbox{\tt {\char'43}{\char'43}} with precedence 6, then this definition would be
illegal:
\bprog
{\small\begin{verbatim}
  a ## b : xs = exp
\end{verbatim}}
\eprog
because \mbox{\tt :} has precedence 5, so the left hand side resolves to
\mbox{\tt (a\ {\char'43}{\char'43}\ x)\ :\ xs}, and this cannot be a pattern binding because \mbox{\tt (a\ {\char'43}{\char'43}\ x)}
is not a valid pattern.

\end{haskellprime}

\outline{
\paragraph*{Translation:}
The general binding form for functions is semantically
equivalent to the equation (i.e.~simple pattern binding):
\[
\mbox{$\it x\ \makebox{\tt =\ {\char'134}}\ x_1\ \ldots \ x_k\ \makebox{\tt ->\ case\ (}x_1\makebox{\tt ,\ }\ldots \makebox{\tt ,\ }x_k\makebox{\tt )\ of}$}
\ba[t]{lcl}
\mbox{$\it \makebox{\tt (}p_{11},\ \ldots ,\ p_{1k}\makebox{\tt )}\ match_1$}  \\
\mbox{$\it \ldots $} \\
\mbox{$\it \makebox{\tt (}p_{n1},\ \ldots ,\ p_{nk}\makebox{\tt )}\ match_n$}
\ea\]
where the \mbox{$\it x_i$} are new identifiers.
}

\subsubsection{Pattern bindings}
\label{patbind}
\index{pattern binding}

A pattern binding binds variables to values.  A {\em simple} pattern
binding has form \mbox{$\it p\ =\ e$}.
\index{simple pattern binding}
The pattern \mbox{$\it p$} is
matched ``lazily'' as an irrefutable pattern\index{irrefutable
pattern}, as if there were an implicit \mbox{\tt {\char'176}} in front 
of it.  See the translation in
Section~\ref{let-expressions}.

The {\em general} form of a pattern binding is \mbox{$\it p\ match$}, where a
\mbox{$\it match$} is the same structure as for function bindings above; in other
words, a pattern binding is:
\[\ba{rcl}
\mbox{$\it p$} & \mbox{$\it \makebox{\tt |}\ \hprime{gs_{1}}$}   & \mbox{$\it \makebox{\tt =}\ e_{1}$} \\
    & \mbox{$\it \makebox{\tt |}\ \hprime{gs_{2}}$}   & \mbox{$\it \makebox{\tt =}\ e_{2}$} \\
    & \mbox{$\it \ldots $} \\
    & \mbox{$\it \makebox{\tt |}\ \hprime{gs_{m}}$}   & \mbox{$\it \makebox{\tt =}\ e_{m}$} \\
    & \multicolumn{2}{l}{\mbox{$\it \makebox{\tt where\ {\char'173}}\ decls\ \makebox{\tt {\char'175}}$}}
\ea\]

%{\em Note}: the simple form
%\WeSay{Yes}
%\mbox{$\it p\ \makebox{\tt =}\ e$} is equivalent to \mbox{$\it p\ \makebox{\tt |\ True\ =}\ e$}.

\outline{
\paragraph*{Translation:}
The pattern binding above is semantically equivalent to this
simple pattern binding:
\[\ba{lcl}
\mbox{$\it p$} &\mbox{\tt =}& \mbox{$\it \makebox{\tt let}\ decls\ \makebox{\tt in}$} \\
    &   & \hprime{\mbox{\tt case\ ()\ of\ }} \\
    &   & \hprime{\mbox{$\it \makebox{\tt \ \ ()\ |\ }gs_{1}\makebox{\tt \ ->\ }e_{1}$}} \\
    &   & \hprime{\mbox{$\it \makebox{\tt \ \ \ \ \ |\ }gs_{2}\makebox{\tt \ ->\ }e_{2}$}} \\
    &   & \hprime{\mbox{$\it \makebox{\tt \ \ \ \ \ \ \ }\ldots $}} \\
    &   & \hprime{\mbox{$\it \makebox{\tt \ \ \ \ \ |\ }gs_{m}\makebox{\tt \ ->\ }e_{m}$}} \\
    &   & \hprime{\mbox{$\it \makebox{\tt \ \ {\char'137}\ ->\ error\ "Unmatched\ pattern"}$}}
\ea\]
}

\section{Static Semantics of Function and Pattern Bindings}
\label{dependencyanalysis}

The static semantics of the function and pattern bindings of
a \mbox{\tt let} expression or \mbox{\tt where} clause
% (including that of the top-level of
% a program that has been translated into a \mbox{\tt let} expression as
% described at the beginning of Section~\ref{modules})
are discussed in this section.

\subsection{Dependency Analysis}
\label{depend-anal}

In general the static semantics are given by \hprime{applying} the
normal Hindley-Milner\index{Hindley-Milner type system} inference
rules.  \hprime{In order to increase polymorphism, these rules are applied to
groups of bindings identified by a \emph{dependency analysis}.}

\begin{haskellprime}

A binding \mbox{$\it b1$} \emph{depends} on a binding \mbox{$\it b2$} in the same list of
declarations if either

\begin{enumerate}
\item \mbox{$\it b1$} contains a free identifier that has no type signature and
is bound by \mbox{$\it b2$}, or
\item \mbox{$\it b1$} depends on a binding that depends on \mbox{$\it b2$}.
\end{enumerate}

A \emph{declaration group} is a minimal set of mutually dependent
bindings. Hindley-Milner type inference is applied to each declaration
group in dependency order. The order of declarations in \mbox{\tt where}/\mbox{\tt let}
constructs is irrelevant.

\end{haskellprime}

% Notes: from http://hackage.haskell.org/trac/haskell-prime/wiki/RelaxedDependencyAnalysis
% 
%     * also tightens up the original wording, which didn't mention that the declarations had to be in the same list and also defined declaration group but not dependency.
%     * defining dependencies between bindings is a little simpler than dependencies between variables.
%     * the dependency analysis transformation formerly listed in this
%     * section is no longer always possible. 

%----------------

\subsection{Generalization}
\label{generalization}

\begin{haskellprime}

The Hindley-Milner type system assigns types to a let-expression in two stages:

\begin{enumerate}
\item The declaration groups are considered in dependency order. For
each group, a type with no universal quantification is inferred for
each variable bound in the group. Then, all type variables that occur
in these types are universally quantified unless they are associated
with bound variables in the type environment; this is called
generalization.

\item Finally, the body of the let-expression is typed.
\end{enumerate}

\end{haskellprime}

% Notes: from http://hackage.haskell.org/trac/haskell-prime/wiki/RelaxedDependencyAnalysis
% 
%     * The original deals with let's consisting of a single binding, instead of declaration groups. Note that we can no longer assume that a let has a single declaration group.
%     * The original does not deal with functions, non-trivial patterns
%     * or recursion. 

For example, consider the declaration
\bprog
{\small\begin{verbatim}
  f x = let g y = (y,y)
        in ...

\end{verbatim}}
\eprog
The type of \mbox{\tt g}'s definition is 
\mbox{$\it a\ \rightarrow\ (a,a)$}.  The generalization step
attributes to \mbox{\tt g} the polymorphic type 
\mbox{$\it \forall\ a.~\ a\ \rightarrow\ (a,a)$},
after which the typing of the ``\mbox{\tt ...}'' part can proceed.

When typing overloaded definitions, all the overloading 
constraints from a single declaration group are collected together, 
to form the context for the type of each variable declared in the group.
For example, in the definition:
\bprog
{\small\begin{verbatim}
  f x = let g1 x y = if x>y then show x else g2 y x
            g2 p q = g1 q p
        in ...
\end{verbatim}}
\eprog
The types of the definitions of \mbox{\tt g1} and \mbox{\tt g2} are both
\mbox{$\it a\ \rightarrow\ a\ \rightarrow\ \makebox{\tt String}$}, and the accumulated constraints are
\mbox{$\it \makebox{\tt Ord}~a$} (arising from the use of \mbox{\tt >}), and \mbox{$\it \makebox{\tt Show}~a$} (arising from the
use of \mbox{\tt show}).
The type variables appearing in this collection of constraints are
called the {\em constrained type variables}.

The generalization step attributes to both \mbox{\tt g1} and \mbox{\tt g2} the type
\[
\mbox{$\it \forall\ a.~(\makebox{\tt Ord}~a,~\makebox{\tt Show}~a)\ \Rightarrow\ a\ \rightarrow\ a\ \rightarrow\ \makebox{\tt String}$}
\]
Notice that \mbox{\tt g2} is overloaded in the same way as \mbox{\tt g1} even though the
occurrences of \mbox{\tt >} and \mbox{\tt show} are in the definition of \mbox{\tt g1}.

If the programmer supplies explicit type signatures for more than one variable
in a declaration group, the contexts of these signatures must be 
identical up to renaming of the type variables.

\subsection{Context Reduction Errors}
\label{context-reduction}
\index{context reduction}
As mentioned in Section~\ref{type-semantics}, the context of a type
may constrain only a type variable, or the application of a type variable
to one or more types.  Hence, types produced by
generalization must be expressed in a form in which all context
constraints have be reduced to this ``head normal form''.
Consider, for example, the
definition:
\bprog
{\small\begin{verbatim}
  f xs y  =  xs == [y]
\end{verbatim}}
\eprog
Its type is given by
\bprog
{\small\begin{verbatim}
  f :: Eq a => [a] -> a -> Bool
\end{verbatim}}
\eprog
and not
\bprog
{\small\begin{verbatim}
  f :: Eq [a] => [a] -> a -> Bool
\end{verbatim}}
\eprog
Even though the equality is taken at the list type, the context must
be simplified, using the instance declaration for \mbox{\tt Eq} on lists,
before generalization.  If no such instance is in scope, a static
error occurs.

Here is an example that shows the need for a
constraint of the form \mbox{$\it C\ (m\ t)$} where m is one of the type
variables being generalized; that is, where the class \mbox{$\it C$} applies to a type
expression that is not a type variable or a type constructor.
Consider:
\bprog
{\small\begin{verbatim}
  f :: (Monad m, Eq (m a)) => a -> m a -> Bool
  f x y = return x == y
\end{verbatim}}
\eprog
The type of \mbox{\tt return} is \mbox{\tt Monad\ m\ =>\ a\ ->\ m\ a}; the type of \mbox{\tt (==)} is
\mbox{\tt Eq\ a\ =>\ a\ ->\ a\ ->\ Bool}.  The type of \mbox{\tt f} should be
therefore \mbox{\tt (Monad\ m,\ Eq\ (m\ a))\ =>\ a\ ->\ m\ a\ ->\ Bool}, and the context
cannot be simplified further.

The instance declaration derived from a data type \mbox{\tt deriving} clause
(see Section~\ref{derived-decls})
must, like any instance declaration, have a {\em simple} context; that is,
all the constraints must be of the form \mbox{$\it C\ a$}, where \mbox{$\it a$} is a type variable.
For example, in the type
\bprog
{\small\begin{verbatim}
  data Apply a b = App (a b)  deriving Show
\end{verbatim}}
\eprog
the derived Show instance will produce a context \mbox{\tt Show\ (a\ b)}, which
cannot be reduced and is not simple; thus a static error results.

\subsection{Monomorphism}
\label{monomorphism}
Sometimes it is not possible to generalize over all the type variables
used in the type of the definition.
For example, consider the declaration\nopagebreak[4]
\bprog
{\small\begin{verbatim}
  f x = let g y z = ([x,y], z)
        in ...
\end{verbatim}}
\eprog
In an environment where \mbox{\tt x} has type \mbox{$\it a$},
the type of \mbox{\tt g}'s definition is 
\mbox{$\it a\ \rightarrow\ b\ \rightarrow\ \makebox{\tt ([}a\makebox{\tt ]},b\makebox{\tt )}$}.
The generalization step attributes to \mbox{\tt g} the type 
\mbox{$\it \forall\ b.~\ a\ \rightarrow\ b\ \rightarrow\ \makebox{\tt ([}a\makebox{\tt ]},b\makebox{\tt )}$};
only \mbox{$\it b$} can be universally quantified because \mbox{$\it a$} occurs in the
type environment.
We say that the type of \mbox{\tt g} is {\em monomorphic in the type variable \mbox{$\it a$}}.
\index{monomorphic type variable}

The effect of such monomorphism is that the first argument of all 
applications of \mbox{\tt g} must be of a single type.  
For example, it would be valid for
the ``\mbox{\tt ...}'' to be
\bprog
{\small\begin{verbatim}
  (g True, g False)
\end{verbatim}}
\eprog
(which would, incidentally, force \mbox{\tt x} to have type \mbox{\tt Bool}) but invalid
for it to be 
\bprog
{\small\begin{verbatim}
  (g True, g 'c')
\end{verbatim}}
\eprog
In general, a type \mbox{$\it \forall\ \overline{u}.~cx\ \Rightarrow\ t$}
is said to be {\em monomorphic}
\index{monomorphic type variable}
in the type variable \mbox{$\it a$} if \mbox{$\it a$} is free in
\mbox{$\it \forall\ \overline{u}.~cx\ \Rightarrow\ t$}.

It is worth noting that the explicit type signatures provided by \Haskell{}
are not powerful enough to express types that include monomorphic type
variables.  For example, we cannot write
\bprog
{\small\begin{verbatim}
  f x = let 
          g :: a -> b -> ([a],b)
          g y z = ([x,y], z)
        in ...
\end{verbatim}}
\eprog
because that would claim that \mbox{\tt g} was polymorphic in both \mbox{\tt a} and \mbox{\tt b}
(Section~\ref{type-signatures}).  In this program, \mbox{\tt g} can only be given
a type signature if its first argument is restricted to a type not involving
type variables; for example
\bprog
{\small\begin{verbatim}
  g :: Int -> b -> ([Int],b)
\end{verbatim}}
\eprog
This signature would also cause \mbox{\tt x} to have type \mbox{\tt Int}.

\subsection{The Monomorphism Restriction}
\index{monomorphism restriction}
\label{sect:monomorphism-restriction}

\Haskell{} places certain extra restrictions on the generalization
step, beyond the standard Hindley-Milner restriction described above,
which further reduces polymorphism in particular cases.

The monomorphism restriction depends on the binding syntax of a
variable.  Recall that a variable is bound by either a {\em function
binding} or a {\em pattern binding}, and that a {\em simple} pattern
binding is a pattern binding in which the pattern consists of only a
single variable (Section~\ref{pattern-bindings}).

The following two rules define the monomorphism restriction:
\outline{
\paragraph*{The monomorphism restriction}

\begin{description}
\item[Rule 1.]
We say that a given declaration group is {\em unrestricted} if and only if:
\begin{description}
\item[(a):]
every variable in the group is bound by a function binding or a simple
pattern binding\index{simple pattern binding} (Section~\ref{patbind}), {\em and}
\item[(b):]
an explicit type signature is given for every variable in the group
that is bound by simple pattern binding.
\end{description}
The usual Hindley-Milner restriction on polymorphism is that
only type variables that do not occur free in the environment may be generalized.
In addition, {\em the constrained type variables of
a restricted declaration group may not be generalized}
in the generalization step for that group.
(Recall that a type variable is constrained if it must belong
to some type class; see Section~\ref{generalization}.)
% 
% \item[Rule 1.]
% The variables of a given declaration group are monomorphic in
% all their constrained type variables if and only if:
% \begin{description}
% \item[either (a):]
% one or more variables in the declaration group 
% is bound by a non-simple pattern binding.
% \item[or (b):]
% one or more variables in the declaration group is bound 
% by a simple pattern binding, and
% no type signature is given for any of the variables in the group.
% \end{description}

\item[Rule 2.]
Any monomorphic type variables that remain when type inference for
an entire module is complete, are considered {\em ambiguous}\index{ambiguous type},
and are resolved to particular types using the defaulting 
rules (Section~\ref{default-decls}).
\end{description}
}

% When all variables in a declaration group are declared using function
% binding the monomorphism restriction will not apply.  Any variable
% declared in a non-simple pattern binding will invoke monomorphism for
% the entire group containing it.  Simple pattern bindings will be
% monomorphic unless a type signature is supplied.
%

\paragraph*{Motivation}

Rule 1 is required for two reasons, both of which are fairly subtle.
\begin{itemize}
\item
{\em Rule 1 prevents computations from being unexpectedly repeated.}
For example, \mbox{\tt genericLength} is a standard function (in library \mbox{\tt Data.List}) whose
type is given by
\bprog
{\small\begin{verbatim}
  genericLength :: Num a => [b] -> a
\end{verbatim}}
\eprog
Now consider the following expression:
\bprog
{\small\begin{verbatim}
  let { len = genericLength xs } in (len, len)
\end{verbatim}}
\eprog
It looks as if \mbox{\tt len} should be computed only once, but without Rule~1 it might
be computed twice, once at each of two different overloadings.  If the 
programmer does actually wish the computation to be repeated, an explicit
type signature may be added:
\bprog
{\small\begin{verbatim}
  let { len :: Num a => a; len = genericLength xs } in (len, len)
\end{verbatim}}
\eprog

\item {\em Rule~1 prevents ambiguity.}  For example, consider the declaration
group
\bprog
{\small\begin{verbatim}
  [(n,s)] = reads t
\end{verbatim}}
\eprog
Recall that \mbox{\tt reads} is a standard function whose type is given by the
signature
\bprog
{\small\begin{verbatim}
  reads :: (Read a) => String -> [(a,String)]
\end{verbatim}}
\eprog
Without Rule~1, \mbox{\tt n} would be assigned the 
type \mbox{$\it \forall\ a.~\makebox{\tt Read}~a\ \Rightarrow\ a$} 
and \mbox{\tt s} the type \mbox{$\it \forall\ a.$} \mbox{$\it \makebox{\tt Read}~a$} \mbox{$\it \Rightarrow\ \makebox{\tt String}$}.
The latter is an invalid type, because it is inherently ambiguous.
It is not possible to determine at what overloading to use \mbox{\tt s}, nor
can this be solved by adding a type signature for \mbox{\tt s}.
Hence, when {\em non-simple} pattern bindings\index{simple pattern binding}
are used (Section~\ref{patbind}), the types inferred are 
always monomorphic in their constrained type variables, irrespective of whether
a type signature is provided.
In this case, both \mbox{\tt n} and \mbox{\tt s} are monomorphic in \mbox{$\it a$}.

The same constraint applies to pattern-bound functions.  For example, in
\bprog
{\small\begin{verbatim}
  (f,g) = ((+),(-))
\end{verbatim}}
\eprog
both \mbox{\tt f} and \mbox{\tt g} are monomorphic regardless of any type
signatures supplied for \mbox{\tt f} or \mbox{\tt g}.
\end{itemize}

Rule~2 is required because there is no way to enforce monomorphic use
of an {\em exported} binding, except by performing type inference on modules
outside the current module.  Rule~2 states that the exact types of all
the variables bound in a module must be determined by that module alone, and not
by any modules that import it.
\bprog
{\small\begin{verbatim}
  module M1(len1) where
    default( Int, Double )
    len1 = genericLength "Hello"

  module M2 where
    import M1(len1)
    len2 = (2*len1) :: Rational
\end{verbatim}}
\eprog
When type inference on module \mbox{\tt M1} is complete, \mbox{\tt len1} has the 
monomorphic type \mbox{\tt Num\ a\ =>\ a} (by Rule 1).  Rule 2 now states that
the monomorphic type variable \mbox{\tt a} is ambiguous, and must be resolved using
the defaulting rules of Section~\ref{default-decls}.
Hence, \mbox{\tt len1} gets type \mbox{\tt Int}, and its use in \mbox{\tt len2} is type-incorrect.
(If the above code is actually what is wanted, a type signature on
\mbox{\tt len1} would solve the problem.)

This issue does not arise for nested bindings, because their entire scope is 
visible to the compiler.

\paragraph*{Consequences}

The monomorphism rule has a number of consequences for the programmer.
Anything defined with function syntax usually
generalizes as a function is expected to.  Thus in
\bprog
{\small\begin{verbatim}
  f x y = x+y
\end{verbatim}}
\eprog
the function \mbox{\tt f} may be used at any overloading in class \mbox{\tt Num}.
There is no danger of recomputation here.  However, the same function
defined with pattern syntax:
\bprog
{\small\begin{verbatim}
  f = \x -> \y -> x+y
\end{verbatim}}
\eprog
requires a type signature if \mbox{\tt f} is to be fully overloaded.
Many functions are most naturally defined using simple pattern
bindings; the user must be careful to affix these with type signatures
to retain full overloading.  The standard prelude contains many
examples of this:
\bprog
{\small\begin{verbatim}
  sum  :: (Num a) => [a] -> a
  sum  =  foldl (+) 0  
\end{verbatim}}
\eprog

Rule~1 applies to both top-level and nested definitions.  Consider
\bprog
{\small\begin{verbatim}
  module M where
    len1 = genericLength "Hello"
    len2 = (2*len1) :: Rational
\end{verbatim}}
\eprog
Here, type inference finds that \mbox{\tt len1} has the monomorphic type (\mbox{\tt Num\ a\ =>\ a});
and the type variable \mbox{\tt a} is resolved to \mbox{\tt Rational} when performing type
inference on \mbox{\tt len2}.

% Even when a function is defined using a function binding, it may still
% be made monomorphic by another variable in the same declaration group.
% Since groups defined through mutually recursive functions need not be
% syntacticly adjacent, it may be difficult to see where overloading is
% being lost.  In this example \mbox{\tt fact'} is defined with a pattern binding
% and forces \mbox{\tt fact} to be monomorphic in the absence of a type signature
% for either \mbox{\tt fact} or \mbox{\tt fact'}.  This would in turn result in an error as
% per Rule~2.
% \bprog
% 
% module Mod1(fact)
% import Mod2
% fact 0 = 1
% fact n = n*fact'(n-1)
% 
% module Mod2(fact')
% import Mod1
% fact' = fact
% 
% \eprog

\section{Kind Inference}
\index{kind}
\index{kind inference}
\label{kindinference}

This section describes the rules that are used to perform {\em kind
inference}, i.e. to calculate a suitable kind for each type
constructor and class appearing in a given
program.

The first step in the kind inference process is to arrange the set of
datatype, synonym, and class definitions into dependency groups.  This can
be achieved in much the same way as the dependency analysis for value
declarations that was described in Section~\ref{dependencyanalysis}.
For example, the following program fragment includes the definition
of a datatype constructor \mbox{\tt D}, a synonym \mbox{\tt S} and a class \mbox{\tt C}, all of
which would be included in the same dependency group:
\bprog
{\small\begin{verbatim}
  data C a => D a = Foo (S a)
  type S a = [D a]
  class C a where
      bar :: a -> D a -> Bool
\end{verbatim}}
\eprog
The kinds of variables, constructors, and classes within each group
are determined using standard techniques of type inference and
kind-preserving unification \cite{jones:cclasses}.  For example, in the
definitions above, the parameter \mbox{\tt a} appears as an argument of the
function constructor \mbox{\tt (->)} in the type of \mbox{\tt bar} and hence must
have kind $\ast$.  It follows that both \mbox{\tt D} and \mbox{\tt S} must have
kind $\ast\rightarrow\ast$ and that every instance of class \mbox{\tt C} must
have kind $\ast$.

% The kind of any constructor that is exported from a module without
% exposing its definition must be specified explicitly in the interface
% for that module.  This is described in Section~\ref{somebodyneedstofillthisin}.

It is possible that some parts of an inferred kind may not be fully
determined by the corresponding definitions; in such cases, a default
of $\ast$ is assumed.  For example, we could assume an arbitrary kind
$\kappa$ for the \mbox{\tt a} parameter in each of the following examples:
\bprog
{\small\begin{verbatim}
  data App f a = A (f a)
  data Tree a  = Leaf | Fork (Tree a) (Tree a)
\end{verbatim}}
\eprog
This would give kinds
$(\kappa\rightarrow\ast)\rightarrow\kappa\rightarrow\ast$ and
$\kappa\rightarrow\ast$ for \mbox{\tt App} and \mbox{\tt Tree}, respectively, for any
kind $\kappa$, and would require an extension to allow polymorphic
kinds.  Instead, using the default binding $\kappa=\ast$, the
actual kinds for these two constructors are
$(\ast\rightarrow\ast)\rightarrow\ast\rightarrow\ast$ and
$\ast\rightarrow\ast$, respectively.

Defaults are applied to each dependency group without consideration of
the ways in which particular type constructor constants or classes are
used in later dependency groups or elsewhere in the program.  For example,
adding the following definition to those above does not influence the
kind inferred for \mbox{\tt Tree} (by changing it to
$(\ast\rightarrow\ast)\rightarrow\ast$, for instance), and instead
generates a static error because the kind of \mbox{\tt []}, $\ast\rightarrow\ast$,
does not match the kind $\ast$ that is expected for an argument of \mbox{\tt Tree}:
\bprog
{\small\begin{verbatim}
  type FunnyTree = Tree []     -- invalid
\end{verbatim}}
\eprog
This is important because it ensures that each constructor and class are
used consistently with the same kind whenever they are in scope.

%**~footer

% Local Variables: 
% mode: latex
% End:
